\documentclass[12pt]{article}  
\usepackage{amsmath}
\usepackage{amssymb}

\usepackage{booktabs}
\usepackage{array}

\usepackage{hyperref}
\providecommand{\href}[1]{\url{#1}}
\usepackage{iftex}
\usepackage[utf8]{inputenc}  
\usepackage{graphicx}
\setlength{\parindent}{0pt}
\setlength{\parskip}{6pt plus 2pt minus 1pt}
\usepackage[normalem]{ulem}
\usepackage[
]{geometry}

\usepackage{amsthm}
\usepackage{float}

\theoremstyle{plain}
\newtheorem{theorem}{Theorem}[section]
\newtheorem{lemma}{Lemma}
\newtheorem{corollary}{Corollary}
\newtheorem{proposition}{Proposition}
\newtheorem{conjecture}{Conjecture}

\theoremstyle{remark}
\newtheorem{remark}{Remark}
\newtheorem{note}{Note}
\newtheorem{claim}{Claim}

\theoremstyle{definition}
\newtheorem{definition}{Definition}[section]  
\newtheorem{condition}{Condition}
\newtheorem{problem}{Problem}[section]
\newtheorem{example}{Example}[section]

\renewcommand\qedsymbol{$\blacksquare$}  

\usepackage{mathtools}
\usepackage{listings}

\usepackage{framed}

\ProvidesPackage{iftex}

\ifxetex

\usepackage{unicode-math}
% \setmainfont{STIX Two Text}
% \setmathfont{STIX Two Math}
% \setmathfont[range={\mathcal,\mathbfcal},StylisticSet=1]{STIX Two Math}

% \setmonofont{JetBrains Mono}

\fi

\providecommand{\tightlist}{\setlength{\itemsep}{0pt}\setlength{\parskip}{0pt}} 

\usepackage{csquotes}
\title{Limits}
\author{ICPR}
\date{January 25, 2021}

\begin{document}


\maketitle

{
\setcounter{tocdepth}{3}
\tableofcontents
}

\section{Proving limits with arbitrary
parameters}

Given \(L\) as a constant, suppose

\[\lim_{x \rightarrow 0}{f(x)} = L\]

Prove

\[\lim_{x \rightarrow 0}{f(2x)} = L\]

This means we suppose

\[\forall\epsilon_{2} > 0,\ \exists\delta_{2} > 0,\ 0 < \left| x_{2} - 0 \right| < \delta_{2} \Rightarrow \left| f\left( x_{2} \right) - L \right| < \epsilon_{2}\]

We need to prove

\[\forall\epsilon > 0,\ \exists\delta > 0,\ 0 < |x - 0| < \delta \Rightarrow \left| f(2x) - L \right| < \epsilon\]


\subsection{Proving the above
statement}

Let \(\epsilon > 0\)

Show the existence of delta to make the statement hold. Unfortunately,
you can't find the delta. However, we can pick a particular value of the
quantified variable labeled under for every. Choose one that is
advantageous to the proof.

Pick \(\epsilon_{2} = \epsilon\). Then

\[\exists\delta_{2} > 0,\ \forall x_{2}\ if\ 0 < \left| x_{2} - 0 \right| < \delta_{2} \Rightarrow \ \left| f\left( x_{2} \right) - L \right| < \epsilon_{2}\]

Now, \(\epsilon_{2}\) must always be picked prior to delta because the
existence of delta depends on epsilon.

Pick \(\delta = \_\_\_\) (leave that blank for now).

Assume \(0 < |x - 0| < \delta\)

Prove \(\left| f(2x) - L \right| < \epsilon\)

Consider (pick) \(x_{2} = 2x\). Why are we allowed to do this? Because
we allowed the blue statement to be true, where \(\forall x_{2}\)
appears somewhere in the blue statement. Why did we choose \(x_{2}\) to
be that? Because it is advantageous to solving the problem. Now, the
statement revisited:

\[\exists\delta_{2} > 0,\ \forall x_{2}\ if\ 0 < \left| x_{2} - 0 \right| < \delta_{2} \Rightarrow \ \left| f\left( x_{2} \right) - L \right| < \epsilon_{2}\]

Because we picked \(x_{2} = 2x\), the red statement above becomes the
statement we need to prove, as
\(\left| f\left( x_{2} \right) - L \right| = |f(2x) - L|\). How do we
prove it? We need to prove the green statement.

When you want to prove an implication, you assume the statement after
the ``if'' and prove the statement after the then.

However, \textbf{when you know an implication is true,} and you want to
use it, you actually prove the statement after the ``if'' and then
conclude the statement after the then. You do not prove the statement
after the ``then'' because we already know the implication is true. This
means \emph{if the green statement is true, the red statement is true.
Proving that the green statement is true verifies the red statement.}

Ask: \(0 < \left| x_{2} - 0 \right| < \delta_{2}\). Can we get this to
hold (what to solve, given \(|x - 0| < \delta\))?

Since (scroll back up to our assumption), we can end up with:

\begin{align*}
|x - 0| &= |x| < \delta \\
0 &< |2x| < 2\delta \\
0 &< \left| x_{2} - 0 \right| < 2\delta
\end{align*}


However, we want to end up with
\(\left| x_{2} - 0 \right| < \delta_{2}\), not \(2\delta\). This is
where we can take advantage of our ``pick \(\delta = \_\_\_\).''
Remember that we are free to pick \(\delta\) to be whatever we want, and
if we pick it to be \(\delta = \frac{\delta_{2}}{2}\), then
\(\delta_{2} = 2\delta\). Now we can substitute:

\[0 < \left| x_{2} - 0 \right| < \delta_{2}\]

Once we know that this (above) is true, we can also conclude (based on
our assumption) that
\(\left| f\left( x_{2} \right) - L \right| < \epsilon_{2}\) is true.

Because of our choices that \(\epsilon_{2} = \epsilon\) and
\(x_{2} = 2x\):

\[\left| f\left( x_{2} \right) - L \right| < \epsilon_{2} \Leftrightarrow \left| f(2x) - L \right| < \epsilon\]

\[\blacksquare\]


\subsubsection{Shortened}

Suppose

\[\forall\epsilon_{2} > 0,\ \exists\delta_{2} > 0,\ 0 < \left| x_{2} - 0 \right| < \delta_{2} \Rightarrow \left| f\left( x_{2} \right) - L \right| < \epsilon_{2}\]

We need to prove

\[\forall\epsilon > 0,\ \exists\delta > 0,\ 0 < |x - 0| < \delta \Rightarrow \left| f(2x) - L \right| < \epsilon\]

Proof. Let \(\epsilon > 0\). Pick \(\epsilon_{2} = \epsilon\), pick
\(\delta = \frac{\delta_{2}}{2}\).

Assume \(|x - 0| < \delta\). Pick \(x_{2} = 2x\). This means

\begin{align*}
2|x - 0| &< 2\delta \\
|2x - 0| &< 2\delta \\
\left| x_{2} - 0 \right| &< 2\delta \\
\left| x_{2} - 0 \right| &< \delta_{2}
\end{align*}


This implies that
\(\left| f\left( x_{2} \right) - L \right| < \epsilon_{2}\) and also
that \(\left| f(2x) - L \right| < \epsilon\).

\(\blacksquare\)


\subsection{Proving the sum of two
limits}

\(\left| f(x) + g(x) - (L + M) \right| < \epsilon\)

\emph{Triangle inequality:} While you can split absolute values over
multiplications, they can't exactly be split when doing addition.
\(|x + y| \leq |x| + |y|\)


\paragraph{Proving the triangle
inequality}

Assume your inequality to be true (but don't write it) (start with the
inequality you want to prove and do some algebra, simplify it, and end
up with something that is necessarily true).

What is \(|a|^{2}\)? It's \(a^{2}\). If \(a\) was negative and you
square it, the negative would go away anyways. Use the square operator
to remove the absolute value.

\begin{align*}
|x + y| &\leq |x| + |y| \\
|x + y|^{2} &\leq \left( |x| + |y| \right)^{2} \\
(x + y)^{2} &\leq \left( |x| + |y| \right)^{2} \\
x^{2} + 2xy + y^{2} &\leq |x|^{2} + 2|x||y| + |y|^{2} \\
x^{2} + 2xy + y^{2} &\leq x^{2} + 2|x||y| + y^{2} \\
2xy &\leq 2|x||y| \\
xy &\leq |xy|
\end{align*}


Is the absolute value of a number always greater than or equal to that
number? Yes. \(xy \leq |xy|\) is true beyond reasonable doubt.

But why does using algebra help prove the statement? Because after each
algebraic step you are making a double-sided implication, where each
step implies and are equivalent to each other.


\subsubsection{Proof}

We assumed the following

\begin{align*}
\forall\epsilon_{2} &> 0,\ \exists\delta_{2} > 0,\ 0 < \left| x_{2} - 0 \right| < \delta \Rightarrow \left| f\left( x_{2} \right) - L \right| < \epsilon_{2} \\
\forall\epsilon_{3} &> 0,\ \exists\delta_{3} > 0,\ 0 < |x_{3} - 0| < \delta \Rightarrow |g(x_{3}) - M| < \epsilon_{3}
\end{align*}


Prove \(\lim_{x \rightarrow 0}{f(x) + g(x)} = L + M\)

Prove \(\forall\epsilon > 0\) if
\(0 < |x - a| < \delta \Rightarrow \left| f(x) + g(x) - (L + M) \right| < \epsilon\)

Let \(\epsilon > 0\) be arbitrary.

Pick \(\epsilon_{2} = \frac{\epsilon}{2}\), (we wouldn't know what
\(\epsilon_{2}\) is, but it is whatever we want any time in the future)
and assume the bottom statement to be true:

\[\exists\delta_{2} > 0\ if\ 0 < \left| x_{2} - a \right| < \delta_{2} \Rightarrow \left| f\left( x_{2} \right) - L \right| < \epsilon_{2}\]

Pick \(\epsilon_{3} = \frac{\epsilon}{2}\), and assume the statement
below to be true:

\[\exists\delta_{3} > 0\ if\ 0 < \left| x_{3} - a \right| < \delta_{3} \Rightarrow \left| g\left( x_{3} \right) - M \right| < \epsilon_{3}\]

Pick
\(\delta = \min\left\{ \delta_{2},\ d_{3} \right\} \Rightarrow \delta \leq \delta_{2},\ \delta \leq d_{3}\)

Prove what we need to solve: assume the ``if'' and prove the ``then.''

Assume \(0 < |x - a| < \delta\)

Prove \(\left| f(x) + g(x) - (L + M) \right| < \epsilon\)

Consider \(x_{2} = x,\ x_{3} = x\).

Since \(0 < |x - a| < \delta\) (we're using the picked blue statement to
justify our substitution):

\begin{align*}
0 \leq |x - a| &= \left| x_{2} - a \right| < \delta \leq \delta_{2} \\
0 < |x - a| &= \left| x_{3} - a \right| < \delta \leq \delta_{3}
\end{align*}


Which helps up conclude:

\begin{align*}
\Rightarrow \left| f\left( x_{2} \right) - L \right| &< \epsilon_{2} \\
\Rightarrow \left| g\left( x_{3} \right) - M \right| &< \epsilon_{3}
 \\
&\left| f(x) + g(x) - (L + M) \right| \\
&= \left| \left( f(x) - L \right) + \left( g(x) - M \right) \right| \\
&\leq \left| f(x) - L \right| + \left| g(x) - M \right| \\
< \epsilon_{2} + \epsilon_{3} &= \epsilon
\end{align*}


In the process, we picked \(\epsilon_{2} = \frac{\epsilon}{2}\) and
\(\epsilon_{3} = \frac{\epsilon}{2}\), as we wanted to end our proof
with epsilon.

We thus conclude \(\left| f(x) + g(x) - (L + M) \right| < \epsilon\).

\[\blacksquare\]


\section{Using existing statements to prove other
statements}

(\(\land\) means ``and''; \(\vee\) means ``or''.)

A statement is \textbf{necessarily true} means we can prove that
statement given using the assumption.

A statement is \textbf{not necessarily true} (considered false, usually
the result of not all \(\forall\) checks passing despite some) when
there is a counterexample.


\subsubsection{Format}

Knowing

\[A \Rightarrow B\]

Prove that

\[C \Rightarrow D\]

\begin{enumerate}
\def\labelenumi{\arabic{enumi}.}
\item
  Assume \(C\) is true.
\item
  Try to prove \(A\) is true, then conclude \(B\) is true for this
  context.
\item
  Use \(B\) to prove \(D\) or find a counterexample.
\end{enumerate}

The end goal is to prove \(D\) is true and \(C\) is false. Remember that
\(C \Rightarrow D\) is logically equivalent to \(\neg C \vee D\).

For example, if we know this is true (statement 1 a.k.a. first equation,
from now on):

\[|x - 2| < 1 \Rightarrow \left| f(x) - 5 \right| < 2\]


\subsubsection{Example 1}

We can prove the following:

\[|x - 2| < 1 \Rightarrow \left| f(x) - 5 \right| < 3\]

Here's the rationale: If we assume \(|x - 2| < 1\), we know that
\(\left| f(x) - 5 \right| < 2\). Because \(2 < 3\),
\(\left| f(x) - 5 \right| < 2 \Rightarrow \left| f(x) - 5 \right| < 3\)
and therefore what we wanted to prove at the start is proven.


\subsubsection{Example 2}

Another statement to prove using the fact that statement 1 is true:

\[|x - 2| < 0.5 \Rightarrow \left| f(x) - 5 \right| < 2\]

By assuming \(|x - 2| < 0.5\), we can conclude \(|x - 2| < 1\), which
immediately implies \(|x - 2| < 0.5\) and
\(\left| f(x) - 5 \right| < 2\) (from the first statement, which turns
out to be what we needed to prove true for the statement we needed to
prove.


\subsubsection{Example 3}

\[|x - 2| < 1 \Rightarrow \left| f(x) - 5 \right| < 1\]

Assume \(|x - 2| < 1\). This means \(\left| f(x) - 5 \right| < 2\), but
can we conclude \(|f(x) - 5| < 1\)? No. Think about it. For example, if
\(3 < f(x) < 7\), then
\(\left| f(x) - 5 \right| < 1 \land \left| f(x) - 5 \right| < 2\), but
if we let \(f(x) = 7.5\), \(|7.5 - 5| = 1.5 < 2\), but \(1,5 \nless 1\).
This is the reasoning why the statement we needed to prove is not
necessarily true. In other words, the steps to prove/disprove a
statement like this is:

\begin{enumerate}
\def\labelenumi{\arabic{enumi}.}
\item
  Assume ``if'' of what to solve
\item
  If the ``if'' of what to solve implies the ``if'' of the first
  equation, then we can guarantee the ``then'' of the first equation is
  true.
\item
  The ``then'' of what to solve must be implied by the ``then'' of the
  first equation. You can also find a counterexample to disprove a
  statement.
\end{enumerate}

A negation of an implication is where the ``if'' is true and the
``then'' is false, always.

Consider \(x = 2\), and \(f(x) = 6.5\). (\(f(x)\) can be anything, as we
did not put any constraints on it.)

\[|x - 2| = 0 < 1 \Rightarrow \left| f(x) - 5 \right| = |6.5 - 5| \nless 1\]

(For what we know that is true must be followed when constructing the
above work.)


\subsubsection{Example 4}

Knowing that:

\[|x - 2| < 1 \Rightarrow \left| f(x) - 5 \right| < 2\]

Prove:

\[|x - 2| < 0.5 \Rightarrow \left| f(x) - 5 \right| < 1\]

\begin{proof}

Assume \(|x - 2| < 0.5\); this implies \(|x - 2| < 1\), and thus implies
\(\left| f(x) - 5 \right| < 2\). If you scroll up, you can find why
\(\left| f(x) - 5 \right| < 1\) is false (not always true). Do this by
giving an example where \(|x - 2| < 0.5\) is true and
\(\left| f(x) - 5 \right| < 1\) is false. Scroll up to question 3 for
when that is used.



\end{proof}
\subsubsection{Example 5}

Knowing that

\[|x - 2| < 1 \Rightarrow \left| f(x) - 5 \right| < 2\]

Prove

\[\exists\delta > 0,\ |x - 2| < \delta \Rightarrow \left| f(x) - 5 \right| < 1\]

To show WTS is false, show

\[\neg\left( \exists\delta > 0,\ |x - 2| < \delta \Rightarrow \left| f(x) - 5 \right| < 1 \right)\]

\[\forall\delta > 0,\ |x - 2| < \delta \land |f(x) - 5| \nless 1\]

(A very small delta will imply that \(x\) can only be 2; any greater
than it will exceed \(\delta\))

Let \(\delta > 0\) be arbitrary. Consider \(x = 2\). This means
\(|x - 2| < \delta\). (Why did we pick \(x = 2\)? Because it is the only
\(x\) that makes \(|x - 2| < \delta\) where delta can be as close to 0
as it wants but never negative)

Consider \(f(x) = 6.5\)

It satisfies \(|x - 2| = 0 < 1 \Rightarrow \left| f(x) - 5 \right| < 2\)

With that checked off, the value we assigned to \(f(x)\) is legal. Then
\(|6.5 - 5| = 1.5 \nless 1\).

This is satisfied, ultimately proving that

\[\forall\delta > 0,\ |x - 2| < \delta \land |f(x) - 5| \nless 1\]

Is true.


\section{Proving that a limit doesn't
exist}

We still have the assumption

\[|x - 2| < 1 \Rightarrow \left| f(x) - 5 \right| < 2\]

We know something about \(f\), but we don't know what the function truly
is. It is just the function that satisfies the assumption.

Now we suppose that

\[\lim_{x \rightarrow 2}{f(x) = L}\]

This means, we know the following is true:

\[\forall\epsilon > 0,\ \exists\delta > 0,\ 0 < |x - 2| < \delta \Rightarrow \left| f(x) - L \right| < \epsilon\]


\subsection{Deciding whether to prove the limit exists or
not}

Can \(L = 2\)? Give an example of a function with that limit being two.
If you think the limit cannot be two, then prove it using the
assumption.

If you say so, it means
\(L - \epsilon < f(x) < L + \epsilon \Rightarrow f(x) = 2\)

Now what do we know about our first assumption? Then
\(\left| f(x) - 5 \right| < 2 \Leftrightarrow 3 < f(x) < 7\)

The two above statements contradict. \(L\) cannot be 2. A.K.A. because
\(f(x) \in (3,\ 7)\) and we wanted \(f(x)\) to output 2, but we know
that can't be the case because \(f(x)\) is restricted.

\emph{To prove an implication is false: A is true, and B is false.}


\subsection{The ``Proving the limit doesn't
exist''}

Prove \(\lim_{x \rightarrow 2}{f(x)} \neq 2\). To prove a limit is not
something, negate the limit definition. So, the limit definition,
negated is:

\[\exists\epsilon > 0,\ \forall\delta > 0,\ \exists x,\ 0 < |x - 2| < \delta \land \left| f(x) - L \right| \geq \epsilon\]

Pick \(\epsilon = 0.5\) (Initially blank and to be written down later.
Could be small as we want.)

Let \(\delta > 0\) be randomly chosen.

Pick \(x = 2 + \min\left\{ \frac{\delta}{2},\frac{1}{2} \right\}\)
(Initially blank. You technically can pick
\(2 + \frac{\delta}{69696969}\) but please don't.)

Prove
\(0 < |x - 2| < \delta \land \left| f(x) - L \right| \geq \epsilon\)

Firstly, prove that the left side is true. Pick a particular \(x\) that
makes the blue statement is true. Remember: our assumption still holds:

\begin{align*}
|x - 2| &< 1 \Rightarrow \left| f(x) - 5 \right| < 2 \\
\Rightarrow 1 &< x < 3 \\
\Rightarrow 3 &< f(x) < 7
\end{align*}


Okay, back to \(0 < |x - 2| < \delta\). Notice that \(x \neq 2\). 2 is
the most obvious answer, but there is a slight technical detail that
needs to be fixed. The statement
\(0 < |x - 2| < \delta \Rightarrow 2 - \delta < x < 2 + \delta\). We
know \(\delta\) can be really small, but not at zero. Therefore, make
\(x\) be slightly larger than 2. This is why we picked
\(x = 2 + \frac{\delta}{2}\), which was picked just now.

Okay, how would you know that \(1 < 2 + \frac{\delta}{2} < 3\)?
\(\delta\) could be anything. The way to solve this is to use the
minimum function. This means
\(x = 2 + \min\left\{ \frac{\delta}{2},\frac{1}{2} \right\}\).

Now since we chose \(x\) to be that:

\begin{align*}
x &\leq 2 + \frac{\delta}{2} < 2 + \delta \\
x &\leq 2 + \frac{1}{2} < 3 \Leftrightarrow x \leq 2.5 < 3
\end{align*}


Compress this, and we have \(x \leq 2 + \frac{\delta}{2}\) and
\(x \leq \frac{5}{2}\).

Since \(\frac{\delta}{2} > 0\) and \(\frac{1}{2} > 0\), \(x > 2\) if we
pick either of them. (\(x\) is designed to be ever so slightly greater
than two.)

\[\Rightarrow 2 - \delta < x < 2 + \delta\]

\[\Leftrightarrow 0 < |x - 2| < \delta\]

What we know: This limit exists. \(\lim_{x \rightarrow 2}{f(x) = L}\).
Can the limit be equal to two? Prove that it can't, which is what we
want to prove.

Back to what we needed to prove:
\(0 < |x - 2| < \delta \land \left| f(x) - L \right| \geq \epsilon\). We
proved \(0 < |x - 2| < \delta\), and now we need to prove
\(\left| f(x) - L \right| \geq \epsilon\).

These are the assumptions we have in our disposal:

\[1 < x < 3 \Rightarrow 3 < f(x) < 7\]

We can confirm that we did declare \(x\) to be slightly greater than two
and we know that it is less than 3 (scroll up), so we can confirm
\(2 < x < 3 \Rightarrow 1 < x < 3\)

Now we finally know \(3 < f(x) < 7\) is true. Knowing this is the fact,
we need to prove \(\left| f(x) - L \right| \geq \epsilon\), which is
\textbf{equivalent} to
\(f(x) \geq L + \epsilon,\ f(x) \leq L - \epsilon\) \textbf{or}
\(f(x) \notin (L - \epsilon,\ L + \epsilon)\) -- if this statement meets
for all intervals of \(f(x)\), \textbf{then we proved what we needed to
solve:} \(\left| f(x) - L \right| \geq \epsilon\) (Recall \(L = 2\))


We didn't choose what epsilon is yet, so can we choose an epsilon in a
way that the interval of \(3 < f(x) < 7\) lies outside of the diagram
above? Yes. Just let \(\epsilon = 0.5\). We know \(L\) is 2, and that
guarantees the interval of \(f(x)\) is outside the interval of
\((L - \epsilon,\ L + \epsilon)\).

Therefore \(f(x)\) does not lie in the interval \((1.5,\ 2.5)\)


A diagram that visualizes it: \(f(x)\) cannot be within the blue zone.

\(\square\)


\subsubsection{Formal Proof, shortened}

Suppose the following are true:

\begin{align*}
|x - 2| &< 1 \Rightarrow \left| f(x) - 5 \right| < 2 \\
&\lim_{x \rightarrow 2}{f(x) = L}\
\end{align*}


Prove \(L \neq 2\).

From the limit presented, we get this, which we can assume holds.

\[\forall\epsilon > 0,\ \exists\delta > 0,\ 0 < |x - 2| < \delta \Rightarrow \left| f(x) - L \right| < \epsilon\]

This is an implication that proves a limit exists. We want to prove that
the limit doesn't exist, so we reverse the epsilon-delta definition of
the limit \(\lim_{x \rightarrow 2}{f(x) = L}\):

\[\exists\epsilon > 0,\ \forall\delta > 0,\ 0 < |x - 2| < \delta \land \left| f(x) - L \right| \geq \epsilon\]

Our task is to prove this with \(L = 2\).

Pick \(\epsilon = 0.5\)

Let \(\delta > 0\) be picked arbitrarily / randomly chosen.

Pick \(x = \min\left\{ 2 + \frac{\delta}{2},\ 2 + \frac{1}{2} \right\}\)

Prove
\(0 < |x - 2| < \delta \land \left| f(x) - L \right| \geq \epsilon\)

We want \(x\) to be a value slightly greater than 2. So, we pick
\(x = \min\left\{ 2 + \frac{\delta}{2},\ 2 + \frac{1}{2} \right\}\) (The
right side of the minimum function prevents \(x\) from being 3 or
greater as it breaks one of our assumptions.

This means \(x \leq 2 + \frac{\delta}{2}\) and
\(x \leq 2 + \frac{1}{2}\). The minimum function also tells us \(x > 2\)
as \(\frac{\delta}{2},\frac{1}{2} > 0\). We can thus write: (also
\(2 + \frac{\delta}{2} < 2 + \delta\))

\begin{align*}
2 &< x < 2 + \delta \\
\Leftrightarrow 0 &< |x - 2| < \delta
\end{align*}


We proved the left side of
\(0 < |x - 2| < \delta \land \left| f(x) - L \right| \geq \epsilon\).
Now to prove \(\left| f(x) - L \right| \geq \epsilon\)

Remember our assumptions:
\(|x - 2| < 1 \Rightarrow \left| f(x) - 5 \right| < 2\)

\[(|x - 2| < 1 \Leftrightarrow 1 < x < 3) \Rightarrow \left( \left| f(x) - 5 \right| < 2 \Leftrightarrow 3 < f(x) < 7 \right)\]

It means \(f(x)\) must be between 3 and 7, non-inclusively. Now, we let
\(L = 2\), remember? This means:

\begin{align*}
\left| f(x) - 2 \right| &\geq \epsilon \\
f(x) &\geq 2 + \epsilon \vee f(x) \leq 2 - \epsilon
\end{align*}


We said we could let \(\epsilon\) be anything, so let's make it 0.5;
this is what happens:

\[f(x) \geq 2.5 \vee f(x) \leq 1.5\]

\[2.5 \leq 3 < f(x) < 7\]

(This says \(f(x)\) must be greater that 2.5 or something else. But
\(f(x)\) has to be between 3 and 7. This means \(f(x)\) must be greater
than 2.5; statement proved -- we proved the other side of the ``and''.)


\subsubsection{Shortened, again}

\begin{align*}
|x - 2| &< 1 \Rightarrow \left| f(x) - 5 \right| < 2 \\
&\lim_{x \rightarrow 2}{f(x) = L}
\end{align*}


\[\exists\epsilon > 0,\ \forall\delta > 0,\ 0 < |x - 2| < \delta \land \left| f(x) - L \right| \geq \epsilon\]

Pick \(\epsilon = 0.5\), let \(\delta > 0\) be picked arbitrarily /
randomly chosen, and pick
\(x = \min\left\{ 2 + \frac{\delta}{2},\ 2 + \frac{1}{2} \right\}\)

Prove \(0 < |x - 2| < \delta\). This statement is the same as
\(2 < x < 2 + \delta\). Because
\(x = \min\left\{ 2 + \frac{\delta}{2},\ 2 + \frac{1}{2} \right\}\),
then \(2 < x\) and \(x < 2 + \delta\). We can thus verify the statement.

Prove \(\left| f(x) - L \right| \geq \epsilon\). Choose
\(\epsilon = \frac{1}{2}\). Then we need to prove that \(f(x)\) is not
in \((L - \epsilon,\ L + \epsilon)\). This means \(f(x)\) is not in
\((2 - 0.5,\ 2 + 2.5)\). We can prove that is the case because \(f(x)\)
is assumed to be in the interval \((3,\ 7)\). Therefore, proven.


\subsection{Preliminary conclusions}

We know \(0 < |x - 2| < \delta\) is true, as \(L\) does not affect it.

\(L = 2.5\)? Choose \(\epsilon = \frac{1}{4}\). Then we need to prove
that \(f(x)\) is not in \((L - \epsilon,\ L + \epsilon)\). This means
\(f(x)\) is not in \((2.25,\ 2 + 2.75)\). We can prove that is the case
because \(f(x)\) is assumed to be in the interval \((3,\ 7)\).
Therefore, proven.

\(L = 8\)? Choose \(\epsilon = \frac{1}{2}\). Then we need to prove that
\(f(x)\) is not in \((L - \epsilon,\ L + \epsilon)\). This means
\(f(x)\) is not in \((8 - 0.5,\ 8 + 0.5)\). We can prove that is the
case because \(f(x)\) is assumed to be in the interval \((3,\ 7)\).
Therefore, proven.

We can conclude preliminary (by realizing) that
\(L \notin (3,\ 7) \Rightarrow\) limit does not exist. However, we have
only checked the values of values out of \(\lbrack 3,\ 7\rbrack\).

\(L = 5\)? This means \(3 < f(x) < 7\) has to hold. Therefore, we think
that it is possible, but it is not a proper answer, as we need to give
an example \(f(x)\) with \(L = 5\), so we can have \(f(x) = 5\): it
satisfies \(3 < f(x) < 7\).

Suppose \(|x - 2| < 1\), prove \(\left| f(x) - 5 \right| < 2\).
\(\left| f(x) - 5 \right| \Rightarrow \left| f(x) - 5 \right| = |5 - 5| < 2\);
therefore, we can be confident that \(L = 5\) is a possibility, but not
completely.

Consider \(L = 4\). Suppose \(|x - 2| < 1\), prove
\(\left| f(x) - 4 \right| < 2\).
\(\left| f(x) - 5 \right| \Rightarrow |4 - 5| = |4 - 5| < 2\);
therefore, we can be confident that \(L = 4\) is a possibility, but not
completely.


\subsection{It is not always what it
seems}

\begin{align*}
|x - 2| &< 1 \Rightarrow \left| f(x) - 5 \right| < 2 \\
&\lim_{x \rightarrow 2}{f(x) = L}
\end{align*}


Remember the definition of a limit at \(x = 2\):

\[\forall\epsilon > 0,\ \exists\delta > 0,\ 0 < |x - 2| < \delta \Rightarrow |f(x) - L| < \epsilon\]

Can \(L = 3\)? We cannot make the conclusion that \(L = 3\) is
impossible, even though making a preliminary conclusion makes it look
like \(L = 3\) is impossible. Suppose \(|x - 2| < 1\), prove
\(\left| f(x) - 5 \right| < 2\).
\(\left| f(x) - 5 \right| \Rightarrow |3 - 5| = |3 - 5| < 2 \Rightarrow 2 < 2\)\ldots{}
however, that does not rule out the possibility of \(L = 3\). If we let
\(f(x)\) be a function represented by the black line with a hole at
\(x = 2\), the graph \(y = f(x)\) approaches 3 from above but doesn't
hit it.


\begin{proof}

Consider \(f(x) = \left\{ \begin{matrix}
(x - 2)^{2} + 3,\ x \neq 2 \\
5,\ x = 2 \\
\end{matrix} \right.\ \). The limit of \(f(x)\) at 2 would be 3, but
\(f(2) = 5\). That point (in red) is irrelevant when we consider limits.
We have a specific condition that \(x \neq 2\). So:

Suppose \(|x - 2| < 1\) (an arbitrary value between 1 and 3). Prove
\(\left| f(x) - 5 \right| < 2\).

Case 1: \(x = 2\). Then \(f(x) = 5\), and
\(\left| f(x) - 5 \right| = |5 - 5| < 2\), sastifying that
\(\left| f(x) - 5 \right| < 2\)

Case 2: \(x \neq 2\). Since \(|x - 2| < 1\)

\begin{align*}
0 &< |x - 2|^{2} < 1^{2} \\
0 &< (x - 2)^{2} < 1 \\
3 &< (x - 2)^{2} + 3 < 4 \\
3 &< f(x) < 4 \\
- 2 &< f(x) - 5 < - 1 \\
1 &< \left| f(x) - 5 \right| < 2
\end{align*}


This proves satisfied that this function works. This means \(L\) can be
3, and we can also come up that \(3 \leq f(x) \leq 7\) (though not
proven yet).



\end{proof}
\section{Limits and Continuity}

Sided-limits are only meant for 2-dimensional visualizations and can be
used for computations.


\subsection{Two-sided limits}

The \textbf{limit} at a point \(a\) of the function \(f(x)\) is equal to
the number \(L\).

\begin{enumerate}
\def\labelenumi{\arabic{enumi}.}
\item
  When you let \(x\) approach \(a\) from the left side \textbf{and} the
  right side of \(a\), the graph of \(f(x)\) gets close to number \(L\).
\item
  The actual value of \(f(a)\) does \textbf{NOT} matter at all.
\end{enumerate}

\[\lim_{x \rightarrow a}{f(x) = L}\]


\subsection{One-sided limits}

A limit from the left side is
\(\lim_{x \rightarrow a^{-}}{f(x) = L_{1}}\). A limit from the right
side is \(\lim_{x \rightarrow a^{+}}{f(x) = L_{2}}\). If the two values
match, the limit of \(f(x)\) at \(a\) from both sides exist.

The formal definition of a left side limit is if \(0 < a - x < \delta\),
and the right limit is \(a < x < a + \delta\). These two cannot be
generalized for higher dimensions.


\subsection{Continuity}

Intuitively, \(f(x)\) is continuous in an interval if the graph if the
graph is in one piece. \(f(x)\) is continuous at the point \(a\) if the
limit \(a\) is equal to \(f(x)\):

\[\lim_{x \rightarrow a}{f(x) = f(a)}\]

We can only define that a function is continuous on a point, but using
that we can:

\(f(x)\) is \textbf{continuous} in an interval means \(f(x)\) is
continuous at every point in the interval. Most functions are
continuous. Transforming a continuous function using any means (such as
+, ×, and composition) does not change whether the transformed interval
is continuous (except division by zero.) Every standard function we know
are continuous in their domains.

E.g.,
\(\lim_{x \rightarrow 1}{\left( x^{2} + 2^{x} \right) = f(1) = 1^{2} + 2^{1} = 3}\)

Find conditions on A and B so that

\[f(x) = \left\{ \begin{matrix}
Ax - B & x \leq 1 \\
3x & 1 < x < 2 \\
Bx^{2} - A & 2 \leq x \\
\end{matrix} \right.\ \]

Is continuous at \(x = 1\) and discontinuous at \(x = 2\).

\textbf{Continuous:} This means \(\lim_{x \rightarrow 1}{f(x) = f(1)}\)
and \(\lim_{x \rightarrow 2}{f(x) \neq f(2)}\). Fortunately, all of the
pieces of the piecewise functions are continuous. To evaluate this
question, we need to use left and right sided limits.

\[\lim_{x \rightarrow 1^{-}}{f(x)} = \lim_{x \rightarrow 1^{-}}{Ax - B} = A - B\]

\[\lim_{x \rightarrow 1^{+}}{f(x)} = \lim_{x \rightarrow 1^{+}}{3x} = 3x = 3\]

\[A - B = 3 = \lim_{x \rightarrow 1}{f(x) = f(1)}\]

As long \(A - B = 3\), then \(f(x)\) is continuous at 1.

\textbf{Discontinuous:} This means
\(\lim_{x \rightarrow 2}{f(x) \neq 2}\)

\[\lim_{x \rightarrow 2^{-}}{f(x) =}\lim_{x \rightarrow 2^{-}}3x = 6\]

\[\lim_{x \rightarrow 2^{+}}{f(x) = \lim_{x \rightarrow 2^{+}}(Bx^{2} - A)} = 4B - A\]

Now, we need to make sure \(\lim_{x \rightarrow 2}{f(x)} \neq f(2)\). If
this two-sided limit were to exist, the left limit and the right limit
would match, meaning
\(6 = 4B - A = \lim_{x \rightarrow 2}{f(x)} = f(2)\). When would the
limit not exist? \(6 \neq 4B - A\). Solve it.

\[\left\{ \begin{matrix}
6 \neq 4B - A \\
3 = A - B \\
\end{matrix} \right.\ \]


\subsection{Limits approaching
infinity}

We allow \(a\) to be infinite as well. This is understood intuitively
as:

As \(x\) becomes larger and largher, \(f(x)\) gets closer to a number
\(L\).

\[\lim_{x \rightarrow \infty}{f(x) = L}\]

We can also have \(a\) being minus infinity:

\(\lim_{x \rightarrow - \infty}{f(x) = L}\)


\subsection{Limit Properties}

First, we \textbf{need} to know the following limits to exist:

\[\lim_{x \rightarrow a}{f(x) = L},\lim_{x \rightarrow a}{g(x) = M}\]

If we are not certain these limits exist, we cannot use the properties
below.

Then we can attain

\[\lim_{x \rightarrow a}{\left\lbrack f(x) + g(x) \right\rbrack = L + M}\]

\[\lim_{x \rightarrow a}{\left\lbrack f(x) - g(x) \right\rbrack = L - M}\]

\[\lim_{x \rightarrow a}{c \cdot f(x) = c \cdot L}\]

\[\lim_{x \rightarrow a}{\lbrack f(x) \cdot g(x)\rbrack = L \cdot M}\]

\[\lim_{x \rightarrow a}{\frac{f(x)}{g(x)} = \frac{L}{M}},\ M \neq 0\]

Third property onward, proving these become extremely hard.

If \(f(x)\) is continuous, we can move the limit \textbf{inside} \(f\):

\[\lim_{x \rightarrow a}{f\left( g(x) \right) = f\left\lbrack \lim_{x \rightarrow a}{g(x)} \right\rbrack}\]


\subsubsection{Limit properties proofs}

\emph{\textbf{If we know two limits is true}}

Firstly, we have this statement brought from the previous subsection:


\begin{equation}
\lim(f)\ \text{and}\lim(g)\ \text{exists} \Rightarrow \lim(f + g) = \lim(f) + \lim(g)\ \tag{1}  \label{eq:1}
\end{equation}


Prove the following statement: Suppose
\(\lim_{x \rightarrow a}{\lbrack f(x) + g(x)\rbrack}\) exists and
\(\lim_{x \rightarrow a}{f(x)}\) exists, then
\(\lim_{x \rightarrow a}{g(x)}\) also exists. Unfortunately, it is not
correct to use (1) at first. You need to prove (1) first, so we know
that it holds, and for (1) to be true we need to know that the limit at
\(f\) and \(g\) exists. However, we do not know for sure that
\(\lim(g)\) exists, meaning we \textbf{cannot assume that (1) is true.}

What would happen if we considered
\(\lim_{x \rightarrow a}{(f + g)} - \lim_{x \rightarrow a}f\). Are we
allowed to consider this? This requires us to use this property:


\begin{equation}
\lim(f)\ \text{and}\lim(g)\ \text{exists} \Rightarrow \lim(f + g)\lim(f - g) = \lim(f) - \lim(g)\ \tag{2}  \label{eq:2}
\end{equation}


This doesn't hold. We do not know that \(\lim_{x \rightarrow a}{g(x)}\)
exists.

What if we consider the function \(f + g\) and the function \(f\)? We
know that \(\lim(f + g)\) exists and so is \(\lim(f)\), so we can deduce
that \(\lim\left( (f + g) - f \right)\).

The key point is to use limit laws, you need to know that the two limits
you would use the laws exist in the first place. We considered
\(\lim(f + g)\) to be one function, and \(\lim(f)\) to be the other,
which is why it is perfectly acceptable to write \(\lim(f + g - f)\). Is
this confusing? Think as if I made the substitution
\(\lim(f + g) = \lim(h)\).

We have \(\lim(f + g - f) = \lim(g)\), \textbf{which we know exists} as
it was derived from a limit that exists. It is precise to say:

\emph{Since} \(\lim(f + g),\lim(f)\) \emph{exists, then we can conclude
that} \(\lim(f + g) - \lim(f)\) \emph{exists, and so on.}


\paragraph{Part 2}

\textbf{T/F?} Suppose
\(\lim_{x \rightarrow a}\left\lbrack f(x) + g(x) \right\rbrack\) exists
but \(\lim_{x \rightarrow a}{f(x)}\) does not exist, then
\(\lim_{x \rightarrow a}{g(x)}\) also does not exist.

\begin{proof} Suppose \(\lim_{x \rightarrow a}{g(x)}\) exists. We know
\(\lim_{x \rightarrow a}\left\lbrack f(x) + g(x) \right\rbrack\) exists.

Consider
\(\lim(f + g) - \lim(g) \Rightarrow \lim(f + g - g) \Rightarrow \lim(f)\)
(we can do this because the limit of \(f + g\) and \(g\) are assumed to
exist), meaning we proved that \(\lim(f)\) exists. However, we have
reached a statement which is contradictory to something we assumed
initially. This means it is not possible for \(\lim(g)\) to exist.

When you reach a contradiction, your ``supposed'' assumption will
automatically falsify itself (or turn true if supposed false). For
regular proof questions, we cannot assume what we want to prove.
However, when working with true/false questions, take a guess and go
both paths.


\paragraph{Part 3}

Suppose \(\lim_{x \rightarrow a}{f(x)}\) exists but
\(\lim_{x \rightarrow a}{g(x)}\) does not exist, then
\(\lim_{x \rightarrow a}{\lbrack f(x) + g(x)\rbrack}\) does not exist.
This statement is true.
\(\lim_{x \rightarrow a}{\lbrack f(x) + g(x)\rbrack}\) cannot exist if
either limit of \(f\) or \(g\) do not exist -- the limit property
\(\lim_{x \rightarrow a}{\lbrack f(x) + g(x)\rbrack}\) only works if the
two functions as individual limits work in the first place. Here, we
know the opposite.



\end{proof}
\subsection{Continuity law for
composition}

Prove the \emph{continuity law for composition:}

If \(f(x)\) is continuous at \(\lim_{x \rightarrow a}{g(x)}\), then we
can move the limit inside:


\begin{equation}
\lim_{x \rightarrow a}{f\left( g(x) \right) = f\left\lbrack \lim_{x \rightarrow a}{g(x)} \right\rbrack}\ \tag{1}  \label{eq:1}
\end{equation}


\textbf{Deducing this implies that composition of continuous functions
is continuous.}

We know that if you put a continuous function inside another continuous
function, you get a composite function. However, we think of continuous
as a graph being in one piece. However, you have no idea.

\begin{framed}

What does it mean for \(f(x)\) to be continuous? It means
\(\lim_{x \rightarrow a}{f(x)} = f(a)\). How would we write this
statement?

\[\forall\epsilon > 0,\ \exists\delta > 0,\ 0 < |x - a| < \delta \Rightarrow |f(x) - f(a)| < \epsilon\]

\end{framed}



Because \(L = f(a)\), we modified our epsilon-delta definition
accordingly. Our assumption is that \(f(x)\) is continuous at
\(\lim_{x \rightarrow a}{g(x)}\), but that should be treated just like a
number. This means, which we know is true:


\begin{equation}
\lim_{x \rightarrow \lim_{x \rightarrow a}{g(x)}}{f(x)} = f\left\lbrack \lim_{x \rightarrow a}{g(x)} \right\rbrack\ \tag{2}  \label{eq:2}
\end{equation}


This is the true definition of if \(f(x)\) is continuous at
\(\lim_{x \rightarrow a}{g(x)}\) (we know
\(\lim_{x \rightarrow a}{g(x)}\) exists, because we must assume
everything in the assumption to be true). This is what the question is
asking.

Now, what does (1) mean? \textbf{WTS:}

\[\forall\epsilon > 0,\ \exists\delta > 0,\ 0 < |x - a| < \delta \Rightarrow \left| f\left( g(x) \right) - f\left\lbrack \lim_{x \rightarrow a}{g(x)} \right\rbrack \right| < \epsilon  \]

This is a modified \(\epsilon - \delta\) definition presented in the box
if you scroll up. This statement is what (1) is trying to say. All we
need to do is prove that (3) is true. This is where we start using our
assumptions. Meaning we know

\[\forall\epsilon_{2} > 0,\exists\delta_{2} > 0,0 < \left| x_{2} - \lim_{x \rightarrow a}{g(x)} \right| < \delta_{2} \Rightarrow \left| f\left( x_{2} \right) - f\left( \lim_{x \rightarrow a}{g(x)} \right) \right| < \epsilon_{2}\  \]

We used assumption (2) to help us construct the limit definition (4)
that we can assume true when proving (3). Our (4) is something that we
use.

\begin{proof} Let \(\epsilon > 0\). Pick \(\epsilon_{2} = \_\_\), which
makes the statement below true.

\[\exists\delta_{2} > 0,\ 0 < \left| x_{2} - \lim_{x \rightarrow a}{g(x)} \right| < \delta_{2} \Rightarrow \left| f\left( x_{2} \right) - f\left( \lim_{x \rightarrow a}{g(x)} \right) \right| < \epsilon_{2}\  \]

Pick \(\delta = \_\_\). It may depend on \(\delta_{2}\).

From (4), suppose \(0 < |x - a| < \delta\)

Prove
\(\left| f\left( g(x) \right) - f\left\lbrack \lim_{x \rightarrow a}{g(x)} \right\rbrack \right| < \epsilon\).

We haven't started the proof yet. We have set up the proof, though.

So, we pick \(\epsilon_{2} = \epsilon\) and \(x_{2} = g(x)\). Then
\(\left| f\left( x_{2} \right) - f\left( \lim_{x \rightarrow a}{g(x)} \right) \right| < \epsilon_{2}\)
from (4) becomes
\(\left| f\left( g(x) \right) - f\left( \lim_{x \rightarrow a}{g(x)} \right) \right| < \epsilon_{2}\)
from (5.0). This means we know this is true:

\[\forall\epsilon_{2} > 0,\exists\delta_{2} > 0,0 < \left| g(x) - \lim_{x \rightarrow a}{g(x)} \right| < \delta_{2} \Rightarrow \left| f\left( g(x) \right) - f\left( \lim_{x \rightarrow a}{g(x)} \right) \right| < \epsilon\  \]

(We are so close! Verify
\(0 < \left| g(x) - \lim_{x \rightarrow a}{g(x)} \right| < \delta_{2}\))

Connection: If
\(a \rightarrow g(a) \approx \lim_{x \rightarrow a}{g(x)} \rightarrow f\left( g(a) \right)\)

If I pick an \(x\) that is very close to \(a\), then
\(x \rightarrow g(x) \rightarrow f\left( g(x) \right)\). We are asking
if \(x\) was near \(a\), can we prove \(f\left( g(x) \right)\) would be
very close to \(f\left( g(a) \right)\)? We then pick an \(\epsilon > 0\)
which gives us an error range around \(f\left( g(a) \right)\), and we
would like to ask: can we find a \(\delta\) range on the left. Then
\(f\left( g(x) \right)\) will be in the epsilon range within
\(f\left( g(a) \right)\)? Then notice that if you start with \(x\), if
you want to get to the right side, it needs to get to \(g(x)\), then
\(g(x)\) would be inputted to \(f\) to obtain \(f\left( g(x) \right)\).
Now, \(\exists\delta_{2}\) around \(g(a)\), then
\(f\left( g(x) \right)\) would be very close to
\(f\left( g(a) \right)\). If we take a look at how \(\delta_{2}\) is set
up, the delta we are looking for is all the way on the left. They are
not even in the same place.


There is an implicit assumption that
\(\lim_{x \rightarrow a}{g(x)} = M\) exists. Then, we know that

\[\forall\epsilon_{3} > 0,\ \exists\delta_{3} > 0,\ 0 < \left| x_{3} - a \right| < \delta_{3} \Rightarrow \left| g\left( x_{3} \right) - \lim_{x \rightarrow a}{g(x)} \right| < \epsilon_{3}\  \]

Pick \(\epsilon_{3} = \delta_{2}\). Then

\[\exists\delta_{3} > 0,\ 0 < \left| x_{3} - a \right| < \delta_{3} \Rightarrow \left| g\left( x_{3} \right) - \lim_{x \rightarrow a}{g(x)} \right| < \epsilon_{3}\]

Pick \(\delta = \delta_{3}\).

Now how can I make sure \(g(x)\) lands in the \(\delta_{2}\) range?

Using the red statement: Give me any \(\epsilon_{3}\) range in the
middle centered around the red point, I will make sure \(g(x)\) lands in
the \(\epsilon_{3}\) error range in the middle by giving a
\(\delta_{3}\) control on the left side. This is why \(\delta_{2}\) is
the control of the \(f\) function, which is the target of the \(g\)
function. This justifies why we picked those variables. These guarantees
\(g(x)\) will be sent to \(f\left( g(x) \right)\). It is not that these
choices are intuitive.

\begin{enumerate}
\def\labelenumi{\arabic{enumi}.}
\item
  Set a target on the right
\item
  Want a control on the left
\item
  Two-leg journey broken up in two pieces
\item
  Have \(g(x)\) land in \(\delta_{2} = \epsilon_{3}\) which grants a
  \(\delta_{3}\) control range on the left
\item
  That will guarantee \(f\left( g(x) \right)\) will land on
  \(\epsilon\).
\end{enumerate}

Considering \(x_{2} = g(x),\ x_{3} = x\)

Since \(0 < |x - a| < \delta = \delta_{3}\)

Then
\(\left| g\left( x_{3} \right) - \lim_{x \rightarrow a}{g(x)} \right| < \epsilon_{3} = \delta_{2} \Rightarrow \left| g(x_{3}) - \lim_{x \rightarrow a}{g(x)} \right| < \epsilon_{3} = \delta_{2}\)

Referring back to equation 6, we can then get:
\(\left| f\left( g(x) \right) - f\left( \lim_{x \rightarrow a}{g(x)} \right) \right| < \epsilon\)



\end{proof}
\section{Limits approaching
infinity}

We do allow \(a\) to be infinite. This can be understood as: As \(x\)
becomes larger and larger, \(f(x)\) becomes closer to \(L\).

\textbf{Vertical Asymptotes:}

\[f(x) = \frac{1}{x - 1}\]

\[\lim_{x \rightarrow 1}{= \pm \infty}\]

\textbf{Horizontal Asymptotes:}

As \(x\) approaches \(\pm \infty\), sometimes \(f(x)\) will attain a
limit. This is the horizontal asymptote. The graph is allowed to cross
the horizontal asymptote. For example, \(\arctan(x)\) gives us
horizontal asymptotes of \(\pm \frac{\pi}{2}\).

Anything divided by infinity, other than infinity, gives zero.

If a limit ends up evaluating to infinity, the limit does not exist.
Also, the sum of a limit that exists and a limit that does not will
exist.


\subsection{Limit Evaluation
Techniques}

ALL INVOLVE GETTING AROUND HOLES. DON'T BREAK THE LIMIT UNTIL YOU'RE
FINISHED! THESE ARE THE EXCEPTIONS TO THE REGULAR RULES.

\textbf{Rational Functions}

Rational functions is a polynomial dividing a polynomial. Factor the top
and bottom and cancel out the factors. Any cases that would result in a
`hole' as you did in Grade 11 will not apply as you are taking limits.

\textbf{Rational functions with} \(\mathbf{x}\) \textbf{approaching
infinity}

Only the term with the highest exponent in the numerator and denominator
matters when computing limits to infinity.

\[\lim_{x \rightarrow \infty}{f(x) = \frac{24x^{6} + 1}{2x^{6} + 30x^{5} + 1} \approx \frac{24x^{6}}{2x^{6}} = 12}\]

To be precise, we divide by the highest exponent on the numerator and
the denominator. In this example, \(x^{6}\) is the highest exponent.

\[\frac{\frac{24x^{6}}{x^{6}} + \frac{1}{x^{6}}}{\frac{2x^{6}}{x^{6}} + \frac{30x^{5}}{x^{6}} + \frac{1}{x^{6}}} = \frac{24 + 0}{2 + 0 + 0} = \frac{24}{2} = 12\]


\[\frac{ax^{m}}{bx^{n}}\]

\(m < n\) indicates the limit is zero.

\(m = n\) indicates the limit is a horizontal asymptote.

\(m > n\) limit is \(\pm \infty\)

\textbf{Square Roots}

Multiply the fraction by its conjugate. Treat the square root as an
entire package. For example:

\[f(x) = \frac{\sqrt{4x + 1} - 3}{x - 2}\]

The square root is \(\sqrt{4x + 1}\). Treat that as one term when
conjugating.

Use this to get around a ``hole''.

\textbf{Small angle approximation}

If the limit is given as \(x \rightarrow 0\) then we can make this
substitution.

\[\sin(x) \approx x \approx \tan(x)\]

This means \(\lim_{x \rightarrow 0}{\frac{\sin(x)}{x} = 1}\).

This works for \(\arcsin(x)\) and \(\arctan(x)\).

Example: \(\lim_{x \rightarrow 0}{\sin(2x)} = 2\).
\(\lim_{x \rightarrow 0}{\frac{\sin^{2}(5x)}{x^{2}} = \frac{\left( \sin(5x) \right)^{2}}{x^{2}}} = \frac{(5x)^{2}}{x^{2}} = \frac{25x^{2}}{x^{2}} = 25\)

Look out for \(\lim_{x \rightarrow 3}\frac{\sin(x - 3)}{x}\).

\textbf{Absolute Values}

Approach from the left and right side separately. The moment you
determine that the expression inside the absolute value is positive or
negative, do what is necessary and remove the absolute value sign. If
the limit exists, they should evaluate to the same values.

For example, approaching from the positive side guarantees \(x\) is
positive, meaning \(|x| = x\); from the negative side guarantees \(x\)
is negative meaning \(|x| = - x\).

It is better to write the actual values of the limits even if it does
not exist, using two single-sided limits.

\textbf{Squeeze Theorem}

The only technique that can be extended to higher dimensions.

To attain \(\lim_{x \rightarrow a}{f(x)}\), you can try to find \(g(x)\)
and \(h(x)\) so that:

\begin{enumerate}
\def\labelenumi{\arabic{enumi}.}
\item
  \(g(x) \leq f(x) \leq h(x)\) near the point \(a\)
\item
  \(\lim_{x \rightarrow a}{g(x) = L = \lim_{x \rightarrow a}{h(x)}}\)
\end{enumerate}

Then we conclude \(\lim_{x \rightarrow a}{f(x) = L}\)

In practice, we want to show \(\lim_{x \rightarrow a}{f(x) = 0}\). Start
at \(|f(x)|\), create a chian of inequality, simplify \(|f(x)|\) to
attain \(|g(x)|\) which has limit 0:

\[\left| f(x) \right| \leq \ldots \leq \left| g(x) \right| \rightarrow 0\ \text{i.e.\ }\lim_{x \rightarrow a}{\left| g(x) \right| = 0}\]

The only possible way for \(\left| f(x) \right|\) to be \(\leq\) and
\(\geq\) to 0 is for \(\lim_{x \rightarrow a}{f(x) = 0}.\)

For example:
\(\lim_{x \rightarrow 0}{x^{2}\sin\left( \frac{1}{x} \right)}\)

Squeeze theorem is used very commonly for
\(\sin\left( \frac{1}{x} \right)\). Remember that
\(\sin(\theta) \in \lbrack - 1,\ 1\rbrack\). This means
\(\left| \sin(\theta) \right| \in \lbrack 0,\ 1\rbrack\). To use squeeze
theorem, take the entire function in absolute value:

\[0 \leq \left| x^{2}\sin\left( \frac{1}{x} \right) \right| \leq \left| x^{2} \right| \cdot 1 \rightarrow 0\]

Thus \(\lim_{x \rightarrow 0}{x^{2}\sin\left( \frac{1}{x} \right)} = 0\)
by squeeze theorem. Refer to the template:
\(\left| f(x) \right| = \left| x^{2}\sin\left( \frac{1}{x} \right) \right|\)
and \(\left| g(x) \right| = \left| x^{2} \right| \cdot 1\)

Note that \(\left| \sin\left( \frac{1}{x} \right) \right| \leq 1\)
because of \(\sin(\theta)\)'s domain.

Squeeze theorem is if you want to show that
\(\lim_{x \rightarrow a}{f(x) = 2}\), you can turn it into
\(\lim_{x \rightarrow a}{f(x) - 2 = 0}\). Regardless of what \(a\) is it
should work regardless of what value it is unless some other problems
independent of squeeze theorem arise.

Discussion: \(y = \sin\left( \frac{1}{x} \right)\). Take a look at
\(y = x^{2}\sin\left( \frac{1}{x} \right)\). Closer to zero, the
function will never be above \(x^{2}\) or \(- x^{2}\). This means
despite the infinite oscillation of \(\sin\left( \frac{1}{x} \right)\),
\(x^{2}\) squeezes the function to zero.

\(x^{2}\) can only be multiplied by a number in the interval
\(\lbrack - 1,\ 1\rbrack\) in
\(y = x^{2}\sin\left( \frac{1}{x} \right)\). This means it can never be
multiplied by anything greater than 1. This is why \(x^{2}\) is the max
value, and \(- x^{2}\) is the minimum value.



\subsection{Types of discontinuity}

Continuous at a point means the limit is the same as the point itself.

\[\lim_{x \rightarrow a}{f(x) = f(a)}\]

To say \(f(x)\) is not continuous, it could fail either case:

\textbf{Removable Discontinuity (HOLE)}

\(\lim_{x \rightarrow a}{f(x) = L}\) exists, but \(f(a) \neq L\) or
\(f(a)\) DNE

In any case, we can redefine the function at \(a\), by setting
\(f(a) = L\).

Removable discontinuity typically arises when a cancellation occurs in a
fraction (rational function):

\[f(x) = \frac{(x - 1)(x + 1)}{x - 1},\ x \neq 1\]

But since the function can be \(x + 1\) if you get rid of the hole, you
can redefine the function as:

\[f_{\text{redefined}}(x) = \left\{ \begin{matrix}
\frac{(x - 1)(x + 1)}{x - 1},\ x \neq 1 \\
2,\ x = 1 \\
\end{matrix} \right.\ \]

What we have done is that despite the hole in \(f(x)\), we can redefine
it to fill in the hole of the function. That is the intuitive thing.

\(f_{\text{redefined}}(x)\) is now continuous at \(a = 1\). Note that
\(f(X) = x + 1\ \)everywhere.

\textbf{Jump Discontinuity}

\[\lim_{x \rightarrow a}{f(x)}\ \text{DNE\ as\ a\ 2-sided\ limit}\]

However, both sided limits exist but are not equal. For example:

\[f(x) = \frac{|x|}{x} = \left\{ \begin{matrix}
1,\ x > 0 \\
 - 1,\ x < 0 \\
\end{matrix} \right.\ \]

This is a true discontinuity as you cannot fix it. We have
\(\lim_{x \rightarrow 0^{-}}{f(x) = - 1}\) and
\(\lim_{x \rightarrow 0^{+}}{f(x) = 1}\)

\textbf{Infinite Discontinuity}

When one or both of the 1-sided limits are \textbf{infinite.}

\[\lim_{x \rightarrow a^{-}}{f(x) = \pm \infty}\ \text{or\ }\lim_{x \rightarrow a^{+}}{f(x) = \pm \infty}\ \text{or\ both}\]

Infinite discontinuity typically arises when a fraction is dividing by
0. However, we require that the numerator to be nonzero. You need to
deduce whether it is \(+ \infty\) or \(- \infty\). This can change
depending on what side you approach from.

Note that \(\frac{0}{0}\) is not allowed. \(\frac{0}{0}\) is called an
indeterminant form, which requires further manipulations. Dividing by
zero and computing limits when doing so can lead to untrue statements, a
case of mathematical fallacy. Sometimes, cancellations may not be
obvious.

\textbf{Essential Discontinuity}

Does not fit the description above. It only comes up in very specific
functions:

\[\lim_{x \rightarrow 0}{\sin\left( \frac{1}{x} \right)}\]

As \(x \rightarrow 0,\ \theta = \frac{1}{x} \rightarrow \infty\). The
oscillations become faster as \(x\) approaches zero. There is no way to
describe the discontinuity. You cannot approach infinite numbers at the
same time, as limits are required to be unique. It cannot be captured by
\(\epsilon - \delta\) limit definitions. The same applies to
\(\lim_{x \rightarrow 0}{\cos\left( \frac{1}{x} \right)}\).


\subsection{The idea of a Derivative}

The derivative of \(f(x)\) from first principles is:

\[f'(a) = \lim_{h \rightarrow 0}\frac{f(a + h) - f(a)}{h}\]

You must do some manipulation to the function to solve this.

Note that \(f'(a)\) is typically different with a different point \(a\).
This can be understood as a function:

\[f'(x) = \lim_{h \rightarrow 0}\frac{f(x + h) - f(x)}{h}\]

Where \(x\) is now thought of as the variable to the function \(f'\).

The three ideas can be presented very differently. They don't seem to be
related at all, but they do end up being the same definition: the first
principles equation. The three definitions are just a coincidence. In
higher dimensions, the three different ideas, simply because they are
very different, they end up being different flavors of derivatives,
hence why they need to be presented.


\subsubsection{Graphically (slope)}

The slope of the line tangent to the graph of the function \(y = f(x)\)
at point \(a\). If you pick a point \(a\) and pick another point, you
can connect them together and that creates a straight line, which you
can compute the slope for it. That new point would be called \(a + h\),
and the corresponding \(y\) values would be \(f(a)\) and \(f(a + h)\).
To find the slope, it would be \(\frac{\Delta y}{\Delta x}\). This is
called the difference quotient which matches the slope of the secant
line. We would like to shift \(a + h\) closer to \(a\), which gives you
a better approximation. The closer you get it, the better the
approximation of the slope would be to the slope of the function at
\(a\). If we approximate that difference to be zero, we get the tangent
line. The tangent line at \(a\) would be defined as the derivative
first-principles equation.

The \textbf{tangent line} at \(x = a\) is defined to be the line with
the \textbf{slope} as the derivative \(f'(a)\) passing through the point
\(\left( a,\ f(a) \right)\).

The point-slope form of a line is
\(y - y_{0} = m\left( x - x_{0} \right)\). Here, the point-slope form of
the line tangent to a point \(a\) of any function would be:

\[y - f(a) = f'(a)(x - a)\]

Where \(y_{0} = f(a),\ m = f'(a),\ x_{0} = a\)


\subsubsection{Sensitivity to Input
changes}

The word ``rate of change'' actually has a more complicated meaning. It
means \emph{the ratio of the changes.} Rate of change is the sensitivity
to input changes. For example, consider a function \(f(x) = 2x\). If you
start with the input \(x = 1\), we have \(f(1) = 2\). If we increase the
input by 1 unit, with \(x = 2\), we have \(f(2) = 4\). The output
changed by two units. If we were to change the input by \(h\) units, the
output is changed by \(2h\) units.

This has nothing to do with the slope. The reason why is because when we
are talking about the slope of a graph, there needs to be a graph in the
first place. However, we usually don't draw graphs all the time. In one
dimension, a graph is a curve on the \(x - y\) axis. In higher
dimensions, it is a surface that floats in mid-air. We don't think of
visual examples here.

If we change the input by a lot, the output will change by a lot. The
meaningful question is not how much the output has change, but rather
the ratio of the changes between the output and the input. It is the
ratio of the change of the output of the function divided by the change
in the input of a function.

\[\frac{\text{change\ in\ output}}{\text{change\ in\ input}}\]

Which is

\[\frac{f(a + h) - f(a)}{h} = \text{Average\ Rate\ of\ Change}\]

It is just a coincidence in one dimension that this idea algebraically
becomes the same expression as the slope of the tangent of the graph,
but the idea is completely different.

If we let \(h \rightarrow 0\), we get the derivative. It measures the
ratio when \textbf{very small} change is done to the input. However, the
ratio of the changes between the output and the input may in fact be
small or large.

(For example: effort becomes less correlated with reception the more
effort you put in.)

The function \(f(x) = x^{2}\), close to zero, changing inputs near it
affects the output a lot less.

Price elasticity of demand has lots of applications here.

The derivative here would be the instantaneous rate of change. The
derivative is a local phenomenon. \(f'(a)\) does not give any
information of \(f(x)\) at other points.


\subsubsection{Velocity}

\(f(t) = s(t)\) is the position of a particle at time \(t\). The average
velocity between times \(t = a\) and \(t = b\) is

\[v_{\text{avg}} = \frac{s(b) - s(a)}{b - a} = \frac{\text{difference\ in\ position}}{\text{time\ used}}\]

This corresponds to the slope of the secant line defined by two points
\(\left( a,\ s(a) \right)\) and \(\left( b,\ s(b) \right)\) on the graph
of \(s(t)\).

The instantaneous velocity, or just simply velocity, at time \(t = a\)
is the derivative:

\[v_{\text{inst}} = v(a) = s'(a) = \lim_{h \rightarrow 0}\frac{s(a + h) - s(a)}{h} = \lim_{t \rightarrow a}\frac{s(t) - s(a)}{t - a}\]

The velocity can again be thought as a function of time:

\[v(t) = s'(t)\]

Velocity also generalizes to higher dimensions. It can be generalized to
be the position of something, having multi-dimensional aspects. This
would be completely different from the rate of change interpretation and
no relation to the slope of the tangent.


\subsection{Calculating derivatives through first
principles}

\[f'(a) = \lim_{h \rightarrow 0}\frac{f(a + h) - f(a)}{h}\]

\[f(x) = x^{3}\]

\begin{align*}
&\frac{(a + h)^{3} - a^{3}}{h} \\
&\frac{a^{3} + 3a^{2}h + 3ah^{2} + h^{3} - a^{3}}{h} \\
&\frac{3a^{2}h + 3ah^{2} + h^{3}}{h} \\
&3a^{2} + 3ah + h^{2} \\
&3a^{2}
\end{align*}


Find \(f(x) = x^{3}\) at (1, 1)

\begin{align*}
y - y_{0} &= m\left( x - x_{0} \right) \\
y - 1 &= m(x - 1) \\
f'(a) &= 3a^{2} = 3a = 3 \\
y - 1 &= 3(x - 1)
\end{align*}


Let \(f'(x) = 0\forall x\) and \(f(x) = 0\). What about \(f(2)\)?
This means the function is not sensitive to input changes meaning
\(f(2) = 0\)

The derivative is a local phenomenon. Let \(f'(x) = 0\) and
\(f(x) = 0\). What about \(f(2)\)? We do not know for sure. For example,
\(f(x) = x^{2},\ f'(x) = 2x\). \(f'(0) = 0\). However, derivative
increases away from zero. At 2, the function becomes more sensitive to
input changes. However, for \(f'(0) = 0\) we only know what happens
to the function at zero and no where else. It Is not about the numerical
difference. As long as \(x\) is not at zero, we do not know what
\(f(x)\) or \(f'(x)\) is. There are infinitely as many values between 0
and 2. It is just as far from zero as 0.2 is, in terms of the number
gap, because all of them are some finite distance away. Knowing
something about the behavior of the function at zero tells us nothing
about the behavior of the function elsewhere.

We however need a more power theorem to prove this statement more
rigorously.

Next:

Let \(s(t) = t^{2} - 6t\). Find the average velocity for \(t\) between
the interval \(\lbrack 4,\ 8\rbrack\) and the instantaneous velocity at
\(t = 8\).

\[v_{\text{avg}} = \frac{s(8) - s(4)}{8 - 4} = \frac{(64 - 48) - (16 - 24)}{4} = \frac{16 + 8}{4} = 6\]

Instantaneous velocity at \(t = 8\)? \(v(t) = s'(t) = 2t - 6\) so, if
you plug in \(v(8) = 16 - 6 = 10\). The velocity is increasing as time
increases. When you take the average of the velocity between 4 and 8,
the number is somewhere in the middle.


\section{Three Important Theorems}

These theorems can't be proven with standard axioms of the real numbers.

The axiom of the least upper bound: TBA


\subsection{Intermediate Value
Theorem}

If \(f(x)\) is \uline{continuous} on \(\lbrack a,\ b\rbrack\) and
\(f(a) < 0,\ f(b) > 0 \Rightarrow \exists c \in (a,\ b)\) with
\(f(c) = 0\) (we mentioned an interval here). \emph{To draw a continuous
curve from below the x-axis to above the x-axis, the curve must cross
the x-axis at least once.}

\textbf{IVT Generalized}

If \(f(x)\) is continuous on \(\lbrack a,\ b\rbrack\) then for any \(K\)
between \(f(a)\), \(f(b)\), there exists \(c \in (a,\ b)\) with
\(f(c) = K\)


\subsection{Extreme Value Theorem
(EVT)}

If \(f(x)\) is continuous on \(\lbrack a,\ b\rbrack\), then \(f(x)\) has
a maximum and minimum on the interval {[}a, b{]}:

\begin{enumerate}
\def\labelenumi{\arabic{enumi}.}
\item
  \(\exists x_{M}\), with \(f\left( x_{M} \right) \geq f(x)\) for all
  \(x \in \lbrack a,\ b\rbrack\) as a maximum.
\item
  \(\mathbf{\exists}x_{m}\), with \(f\left( x_{m} \right) \leq f(x)\)
  for all \(x \in \lbrack a,\ b\rbrack\) as a minimum.
\end{enumerate}

(On a closed interval a continuous curve would have a highest and lowest
point.)


\subsection{Mean value theorem}

If \(f(x)\) is continuous on \(\lbrack a,\ b\rbrack\) and is
differentiable on \((a,\ b)\), then \(\exists c \in (a,\ b)\) with

\[f'(c) = \frac{f(b) - f(a)}{b - a}\]

(In terms of velocity, there is a \(c\) where
\(f'(c) = \text{average\ velocity}\))

\emph{Intuition: if the average speed of a car is
100kmh\textsuperscript{-1}, will there be a point where the car travels
at exactly 100kmh\textsuperscript{-1}?}

These theorems do not help find the exact values -- sometimes, you
cannot find exact values; you must use approximations. It just
guarantees the existence of them.

(Note: differentiable on \(\lbrack a,\ b\rbrack \Rightarrow\)
differentiable on \((a,\ b)\) -- we use the weakest assumption. Rare
situation is when cusps exist on the edge of the interval.)

Jumps and cusps are not allowed. Jumps would be considered
teleportation, and cups are sudden increases in speed with no
acceleration for some time, and that cannot happen in real life.

When the assumption for the theorems is not satisfied, the theorem may
not work. Even if the assumptions are not satisfied, the theorem may
still work, but not always.


\subsubsection{Q}

Show that there are at least 3 solutions in
\(\left( \frac{x}{\pi} \right)^{3} - 2\sin(x) + \frac{1}{2} = 0\)

\(x\) is not solvable. Knowing you can't solve it, use IVT.

\begin{table}[H]
\centering


\begin{tabular}{|m{17em}|m{17em}|}

\hline
\[x\]& \[f(x)\]\\ \hline
\[0\] & \[\frac{1}{2}\] \\ \hline
\[\pi\] & \[\frac{3}{2}\]
\\\hline
\end{tabular}

\end{table}
Interval checking


\section{Squeeze theorem}

For every \(x \neq 0\):

\[- 1 \leq \sin\frac{1}{x} \leq 1\]

\[- x^{2} \leq x^{2}\sin\frac{1}{x} \leq x^{2}\]

A more general picture:

If you have two functions which have the same limit \(L\) at \(a\), and
I have a third function that is squeezed in between them, no matter how
it behaves, it has the same limit.

This means:

Let \(a,\ L\mathbb{\in R}\). Let \(f,\ g,\ h\) be functions defined near
\(a\), except possibly at \(a\)

If:

\begin{itemize}
\item
  \(h(x) \leq g(x) \leq f(x)\)
\item
  \(\lim_{x \rightarrow a}{f(x)} = L\)
\item
  \(\lim_{x \rightarrow a}{h(x)} = L\)
\end{itemize}

Then

\begin{itemize}
\item
  \(\lim_{x \rightarrow a}{g(x)} = L\)
\end{itemize}

Compute \(\lim_{x \rightarrow 0}{x^{2}\sin\frac{1}{x}}\).

\begin{itemize}
\item
  For every \(x \neq 0:\  - x^{2} < x^{2}\sin\frac{1}{x} \leq x^{2}\)
\item
  \(\lim_{x \rightarrow 0}x^{2} = 0\) and
  \(\lim_{x \rightarrow 0}{( - x^{2})} = 0\)
\item
  By the squeeze theorem:
  \(\lim_{x \rightarrow 0}{x^{2}\sin\frac{1}{x} = 0}\)
\end{itemize}


\section{Derivative}

The function is differentiable at a point if the limit exists as the
derivative at the point.

\textbf{Tangent line}

\[y - y_{0} = m\left( x - x_{0} \right)\]

Tangent line at
\(\left( x_{0},\ y_{0} \right) = \left( a,\ f(a) \right)\) has slope
\(m = f'(a)\)

\[e^{\ln(x)} = \ln\left( e^{x} \right)\]

(This works because an inverse function acts as an argument to the
function)

Note: \(\tan'(x) = \sec^{2}(x)\) and
\(\left( \sec(x) \right)' = \sec(x)\tan(x)\)


\subsection{Inverse Function Theorem}

\[\left( f^{- 1} \right)'(x) = \frac{1}{f'\left( f^{- 1}(x) \right)}\]

\begin{align*}
\sin^{- 1}{(x)} &= \frac{1}{\sqrt{1 - x^{2}}} \\
\text{acos}(x) &= - \frac{1}{\sqrt{1 - x^{2}}}
\end{align*}


\[\text{atan}'(x) = \frac{1}{1 + x^{2}}\]

Square roots: \(\sqrt{x} \geq 0\)

\begin{itemize}
\item
  \(x = 4\)
\item
  \(\sqrt{x} = a\) (Square root is defined to simply represent the
  positive side.)
\item
  \(a^{2} = 4 \Rightarrow a = \pm 2\)
\end{itemize}


\subsection{Higher order derivatives}

\[\sin^{2}{(x)} = \sin^{2}(x) + \sin^{2}x^{2}\]

Derivative is

\[2\sin(x)\cos(x) + \cos\left( x^{2} \right)2x + 2\sin\left( x^{2} \right)\cos\left( x^{2} \right)2x\]

Derivative of \(e^{x^{e}}\) is \(e^{x^{e}}ex^{e - 1}\)

\(x^{x}\) (when both the base of the exponent and the exponent is the
same, you cannot take the derivative.) \textbf{The power rule may not be
applied if the exponent is variable}

\[x = e^{\ln(x)} \Rightarrow x^{x} = e^{\left( \ln(x) \right)^{x}} = \left( e^{\ln(x)} \right)^{x} = e^{x \cdot \ln(x)}\]

Using \(x = e^{\ln(x)}\) is workable. The derivative of
\(e^{x \cdot \ln(x)}\) is
\(e^{x \cdot \ln(x)} \cdot \left( 1 \cdot \ln(x) + x\left( \frac{1}{x} \right) \right)\)
(then simplify it using the product rule)


\subsection{Nth order derivatives}

\[f(x) = \ln(x)\]

\begin{align*}
f'(x) &= x^{- 1} \\
f''(x) &= - x^{- 2} \\
f'''(x) &= ( - 1)( - 2)x^{- 3} \\
f^{4}(x) &= ( - 1)( - 2)( - 3)x^{- 4} \\
f^{n}(x) &= ( - 1)^{n - 1}\left( 1 \times 2 \times \cdots \times (n - 1) \right)x^{- 4} = - 1 \cdot (n - 1)! \cdot x^{- 4}
\end{align*}


The \(( - 1)^{n - 1}\) is how we compress signs. For example, the second
derivative ends up\ldots.


Firstly, derivative: \(3x^{2} = A\) at \(x = 1\), meaning \(3 = A\).
This is not enough to guarantee that \(f\) is differentiable at 1.

Second, continuous: \(x^{3} = 3x + b\) at \(x = 1\) means
\(1 = 3 + b \Rightarrow b = - 2\)


\subsection{Exponents and logs}

\[\left( e^{x} \right)' = e^{x}\]

What if base \(e\) are not used?

We will use the fact that \(e^{x}\) and \(\ln(x)\) are inverse
functions.

Use the relation:
\(a^{x} = \left( e^{\ln(a)} \right)^{x} = e^{x \cdot \ln(a)}\)

\[\left( a^{x} \right)' = e^{x \cdot \ln(a)} \cdot \ln(a) = a^{x} \cdot \ln(a)\]


\subsubsection{Natural log:}

\[\left( \ln(x) \right)' = \frac{1}{x}\]

(We define \(\ln(x)\) for which its derivative is exactly
\(\frac{1}{x}\). It does not come from base 10 log or any other base.
\(\ln(x)\) is defined to be the integral of \(\frac{1}{x}\))


\subsubsection{Derivative of logs:}

\[y = \log_{a}{(x) \Leftrightarrow x = a^{y}}\]

Using the relation \(a = e^{\ln(a)}\):

\[x = a^{y} = \left( e^{\ln(a)} \right)^{y} = y \cdot \ln(a)\]

Take the ln of both sides:

\[\ln(x) = \ln\left( e^{y} \cdot \ln(a) \right) = y \cdot \ln(a)\]

Divide by the constant \(\ln(a)\)

\[y = \log_{a}(x) = \frac{\ln(x)}{\ln(a)}\]

(This is simply log change of base. \(\ln(a)\) is simply a vertical
stretch.)

Taking the derivative:
\(y' = \left( \log_{a}(x) \right)' = \frac{1}{\ln(a)} \cdot \frac{1}{x} = \frac{1}{x\ln(a)}\)

I mean:

\begin{align*}
f(x) &= e^{x} \\
f\left( x \cdot \ln(a) \right) &= e^{x \cdot \ln(a)}
\end{align*}



\subsubsection{Logarithmic
differentiation:}

\emph{When rolling a dice and flipping a coin, the probability of
rolling a 1 and getting a head on the coin is the product of the
probabilities of either. With the logical connection AND, the
probabilities multiply together. When you have a lot of things connected
with AND, it is a bunch of things multiplied by each other. The problem
is, the derivative works very well with addition but not very well with
multiplication.}

\begin{align*}
\left( fg \right)' &= f'g + fg' \\
\left( \text{fgh} \right)' &= f'gh + fg'h + fgh'
\end{align*}


(The prime carries from the left to the right.)

We can use ln to change multiplication to addition, division to
subtraction, and exponentiation to multiplication. We have the following
properties:

\begin{align*}
\ln\left( xy \right) &= \ln(x) + \ln(y) \\
\ln\left( \frac{x}{y} \right) &= \ln(x) - \ln(y) \\
\ln\left( x^{y} \right) &= y\ln(x)
\end{align*}


To use it:

\[y(x) = \frac{f(x)^{2}g(x)}{h(x)}
\]

\[\ln\left( y(x) \right) = \ln\left( \frac{f(x)^{2}g(x)}{h(x)} \right)\]

Applying the log rules, we get (note that the square in \(f(x)\) is part
of the function, as \(\ln\left( f(x)^{2} \right) = 2\ln{f(x)}\):

\[\ln\left( y(x) \right) = 2\ln\left( f(x) \right) + \ln\left( g(x) \right) - \ln\left( h(x) \right)\]

Take the derivative against \(x\) on both sides using the chain rule, we
get:

\[\frac{1}{y(x)} \cdot y'(x) = \frac{2}{f(x)} \cdot f'(x) + \frac{1}{g(x)} \cdot g'(x) - \frac{1}{h(x)} \cdot h'(x)\]

Multiply both sides by \(y(x)\)

\begin{align*}
\frac{y(x)}{y(x)} \cdot y'(x) &= 2\frac{y(x)}{f(x)} \cdot f'(x) + \frac{y(x)}{g(x)} \cdot g'(x) - \frac{y(x)}{h(x)} \cdot h'(x) \\
y'(x) &= 2\frac{y(x)}{f(x)} \cdot f'(x) + \frac{y(x)}{g(x)} \cdot g'(x) - \frac{y(x)}{h(x)} \cdot h'(x)
\end{align*}


Okay good luck stuffing \(y(x)\) into this mess

\[y'(x) = 2\frac{\frac{f(x)^{2}g(x)}{h(x)}}{f(x)} \cdot f'(x) + \frac{\frac{f(x)^{2}g(x)}{h(x)}}{g(x)} \cdot g'(x) - \frac{\frac{f(x)^{2}g(x)}{h(x)}}{h(x)} \cdot h'(x)\]


\paragraph{An example}

\[y(x) = \left( \frac{\sqrt{x^{2} + 3}}{\sqrt[3]{x^{2} + 2x}} \right)\]

What is \(\ln\left( y(x) \right)\)? It is

\[\ln\left( y(x) \right) = \frac{1}{2}\ln\left( x^{2} + 3 \right) - \frac{1}{3}\ln\left( x^{2} + 3x \right)\]

\[\frac{1}{y(x)}y'(x) = \frac{1}{2\left( x^{2} + 3 \right)} - \frac{1}{3\left( x^{2} + 3x \right)}\]

\[y'(x) = \frac{y(x)}{2\left( x^{2} + 3 \right)} - \frac{y(x)}{3\left( x^{2} + 3x \right)}\]

\[y'(x) = \frac{\frac{\sqrt{x^{2} + 3}}{\sqrt[3]{x^{2} + 2x}}}{2\left( x^{2} + 3 \right)} - \frac{\frac{\sqrt{x^{2} + 3}}{\sqrt[3]{x^{2} + 2x}}}{3\left( x^{2} + 3x \right)}\]


\section{Inverse functions (computationally
difficult)}

A function starts with a domain and a codomain.

\[f:A \rightarrow B\]


\subsection{One-to-one}

A function is \textbf{one-to-one} when any two different points in the
domain will get sent (by \emph{f}) to two different points in the
codomain. Equivalently:

\textbf{There cannot exist 2 different points that get sent by f to the
same point.}

The simplest logical statement for one-to-one is:

\textbf{f is one-to-one} \(\Leftrightarrow\)
\(\left( f(a) = f(b) \Rightarrow a = b \right)\)

\emph{Why is this the case? If it seems as though we have 2 different
points sent to the same point, then these two points have to be the
same.} What you are concluding is that it is \textbf{not} two different
points.

\[f(a) = f(b) \Rightarrow a = b\]

We feel as though \(a\) and \(b\) are different, but we never said that.
The only time where it seems that two different points are getting sent
to the same point, you didn't start with two different points.

Contrapositive of \(a \neq b \Rightarrow f(a) \neq f(b)\) \emph{There
cannot exist two points that get sent to the same point.}

\textbf{Example:}

\(f(x) = x + 3\). Suppose \(f(a) = f(b)\), \(a + 3 = b + 3\), so
\(a = b \Rightarrow f\) is one-to-one


\subsection{Onto}

The \textbf{image} is all possible values that \(f\) can take (based on
the domain).

\(f\) is \textbf{onto} if: \textbf{image = codomain}

Consider this function, which we can assume to be one-to-one for now:

We have defined our function to be
\(f\mathbb{:R \rightarrow R,\ }f(x) = \arctan(x)\)

Consider \(A\mathbb{= R}\) (on the left) and \(B\mathbb{= R}\) (on the
right). However, as \(x\) can be any real number in
\(A\mathbb{= R,}\arctan(x)\) can take on every value in
\(\left( - \frac{\pi}{2},\frac{\pi}{2} \right)\).

The function is not \textbf{onto,} as we typically define the codomain
for 1D calculus to be \(\mathbb{R}\), but the image does not cover
\(\mathbb{R}\).

However, we can redefine the codomain. If we set
\(f\mathbb{:R \rightarrow}\left( - \frac{\pi}{2},\frac{\pi}{2} \right)\),
then we have made the function onto. Then \(f\) is \textbf{one to one
and onto.}

This means we can remove elements from the codomain by redefining what
the function is supposed to map to in order to make our function onto.

(e.g. in linear algebra,
\(f:\mathbb{R}^{2} \rightarrow \mathbb{R}^{3}\), codomain is
\(\mathbb{R}^{3}\). Yes, your linear transformation can be redundant as
\(\left\lbrack x_{1} + x_{2} + x_{3},\ 0,\ 0 \right\rbrack\) -- if you
redefine the codomain to be smaller you are going to have to get rid of
dimensions of that linear transformation)

\textbf{DOMAIN}

If the function's domain is \(\mathbb{R}\), only consider the point
where the function is defined at that point.


\subsection{Horizontal line test
(one-to-one)}

Given a function \(f:f\) is one-to-one if and only if every possible
horizontal line in the x-y plane intersect the curve at most once.


\subsection{One-to-one and onto}

A function is \textbf{one-to-one and onto}

\[f:A \rightarrow B\]

\textbf{if and only if there exists an inverse function}
\(\mathbf{f}^{\mathbf{- 1}}\)\textbf{:}

\[f^{- 1}:B \rightarrow A\]

Suddenly, \(B\) is the domain.

One-to-one rationale -- if it wasn't, our inverse wouldn't be a function

If it isn't onto, there will be missing correspondences. Better remove
things from the codomain first, because if something is on the codomain
that can't be associated with a point on the domain, then you have a
missing correspondence.

The inverse comes in pairs. The function \(f\) takes \(a \in A\) as
input and gibves \(b \in B\) as output. The function \(f^{- 1}\) takes
\(b \in B\ \) and spits out \(a \in A\) as output of \(f^{- 1}\).

The domain of \(f^{- 1}\) is B and the image of \(f^{- 1}\) is \(A\).

The inverse of \(f^{- 1}\) is \(f\), written as
\(\left( f^{- 1} \right)^{- 1} = f\)

Meaning

\[f^{- 1}\left( f(1) \right) = 1;f\left( f^{- 1}(4) \right) = 4\]

The inverse function and the regular function should cancel each other
out.

There is a one-to-one correspondence between the points in \(A\) and
\(B\).

\begin{align*}
y &= f(x) \\
f^{- 1}(y) &= f^{- 1}\left( f(x) \right) \\
f^{- 1}(y) &= x
\end{align*}


This means if \(y = f(x)\),
\(y = \arctan(x) \Rightarrow f^{- 1}(y) = \tan(y) \Rightarrow x = \tan(y)\)

However:
\(f\mathbb{:R \rightarrow}\left( - \frac{\pi}{2},\frac{\pi}{2} \right)\)

\[f^{- 1}:\left( - \frac{\pi}{2},\frac{\pi}{2} \right)\mathbb{\rightarrow R}\]

The domain is restricted. The domain of the inverse function does not
expand infinitely out. This is how inverse trigonometric functions are
defined.

Although \(f^{- 1}(y)\) is defined to be a formula to \(\tan(y)\), it
has a different domain than the standard \(\tan(y)\) function.

The inverse function only works is because we have restricted the domain
of the input for the inverse. For
\(y \notin \left( - \frac{\pi}{2},\frac{\pi}{2} \right)\), we have
\(\arctan\left( \tan(y) \right) \neq y\). To avoid confusion, instead of
using \(y\), we will be using \(\theta\) for angle and \(x\) for the
real number part.


\subsection{Inverse trig function}

Consider \(f(\theta) = \sin(\theta)\). Let
\(f:A\mathbb{= R \rightarrow}B\mathbb{= R}\)

We have \(\text{image}(f) = \lbrack - 1,\ 1\rbrack\) so we set

\[f:A\mathbb{= R \rightarrow}B\lbrack - 1,\ 1\rbrack\]

We see that for a sine function, it is not 1:1 -- the image is
\(\lbrack - 1,\ 1\rbrack\) so the function is onto. We can fix that by
redefining the codomain is \(\lbrack - 1,\ 1\rbrack\), so
\(\sin(\theta)\) is onto. But \(\sin(\theta)\) is not one-to-one, as it
fails the horizontal line test.

Restrict the domain by removing \(2\pi\) and \(4\pi\) if the circle on
the left is our domain. At least something smaller than all real
numbers. While the original \(f(\theta)\) is the entire sine functiuon,
we will restrict the domain to be the green portion of the curve.

The most common way for \(f(\theta) = \sin(\theta)\) is to choose
\(A = \left\lbrack - \frac{\pi}{2},\frac{\pi}{2} \right\rbrack\). Then
\(f\) is \textbf{one-to-one and onto,} and \(f^{- 1}(x) = \arcsin(x)\)

\[f(\theta) = \sin(\theta):\left\lbrack \frac{-\pi}{2},\frac{\pi}{2} \right\rbrack \rightarrow \lbrack - 1,\ 1\rbrack\]

\[f^{- 1}(x) = \arcsin(x):\lbrack - 1,\ 1\rbrack \rightarrow \left\lbrack - \frac{\pi}{2},\frac{\pi}{2} \right\rbrack\]

Sin takes an angle and returns a ratio; arcsin takes a ratio and returns
an angle.

We restrict around zero for the sake of simplicity.

These restrictions are universally agreed with to give the functions
one-to-one properties. The confusing part about
\(\arcsin\left( \sin(6) \right) \neq 6\) -- how do we deal with angles
that are outside the domain that would make \(\sin(\theta)\) bijective?


\subsection{Working with inverse trigonometric
functions}

Following values:

\[x = \cos\left( \frac{2\pi}{3} \right) = - \frac{1}{2}\]

\[y = \sin\left( \frac{2\pi}{3} \right) = \frac{\sqrt{3}}{2}\]

\(\frac{2\pi}{3}\) is in the restricted domain
\(\lbrack 0,\ \pi\rbrack\). This means
\(\arccos(x) = \arccos\left( \cos\left( \frac{2\pi}{3} \right) \right) = \frac{2\pi}{3}\)

The two parts,
\(\arccos\left( \cos\left( \frac{2\pi}{3} \right) \right)\) are inverses
of each other, giving exactly \(\frac{2\pi}{3}\) back, meaning
\(\arccos\left( - \frac{1}{2} \right) = \frac{2\pi}{3}\). We can only
cancel out inverse functions if the input is in the restriction.

Okay. Now, consider
\(y = \sin\left( \frac{2\pi}{3} \right) = \frac{\sqrt{3}}{2}\). Now what
is
\(\arcsin\left( \sin\left( \frac{2\pi}{3} \right) \right) = \arcsin(y)\)?

Unfortunately,
\(\arcsin\left( \sin\left( \frac{2\pi}{3} \right) \right) \neq \frac{2\pi}{3}\).
However, if you look at how arcsine is defined, it will only take values
between
\(\lbrack - 1,\ 1\rbrack \rightarrow \left\lbrack - \frac{\pi}{2},\frac{\pi}{2} \right\rbrack\).

\begin{itemize}
\item
  We said
  \(y = \sin\left( \frac{2\pi}{3} \right) = \frac{\sqrt{3}}{2}\). Why is
  the representation of
  \(\frac{\sqrt{3}}{2} = \sin\left( \frac{2\pi}{3} \right)\)? Is there
  another way to represent \(\frac{\sqrt{3}}{2}\)? Consider the angle
  \(\frac{\pi}{3}\). Note that
  \(\frac{\sqrt{3}}{2} = \sin\left( \frac{\pi}{3} \right)\)
\end{itemize}

\[\arcsin\left( \sin\left( \frac{2\pi}{3} \right) \right) = \arcsin(y) = \arcsin\left( \sin\left( \frac{\pi}{3} \right) \right)\]

Because \(\frac{\pi}{3}\) falls in the region for how arcsine is
defined, arcsine and sine here act as inverse functions to each other
and give out the same thing. Meaning
\(\arcsin\left( \sin\left( \frac{2\pi}{3} \right) \right) = \arcsin\left( \sin\left( \frac{\pi}{3} \right) \right) = \frac{\pi}{3}\).

\(\frac{\sqrt{3}}{2}\) has multiple ways of being represented. In other
words, if I were to apply arcsine on both sides of the equation, we get
\(\arcsin(y) = \arcsin\left( \sin\left( \frac{2\pi}{3} \right) \right)\)
which is the natural thing to do, but then if you were to do this, sine
and arcsine cannot act as inverse functions as each other, as
\(\frac{2\pi}{3}\) is not in the correct angle domain for sine and
arcsine that would make them inverse functions of each other.

So instead, \(\frac{\sqrt{3}}{2}\) not only can be represented as
\(\sin\left( \frac{2\pi}{3} \right)\), but I could also represent it as
\(\sin\left( \frac{\pi}{3} \right)\).

Now, other cases:
\(x = \cos\left( \frac{2\pi}{3} \right) = - \frac{1}{2}\);
\(\arccos\left( \cos\left( \frac{2\pi}{3} \right) \right) = \frac{2\pi}{3}\)

Now \(\arctan\left( \tan\left( \frac{2\pi}{3} \right) \right)\)

Note that
\(z = \tan\left( \frac{2\pi}{3} \right) = \frac{\sqrt{3}}{- 1}\)

Use the same technique: what between
\(( - \frac{\pi}{2},\frac{\pi}{2})\) if inputted into tan gives us
\(- \sqrt{3}\)?

Use the CAST rule. \(\tan\left( - \frac{\pi}{3} \right) = - \sqrt{3}\).
Using this insight,
\(- \frac{\pi}{3} = \arctan\left( \tan\left( - \frac{\pi}{3} \right) \right)\).
Notice that
\(- \frac{\pi}{3} \in \left( - \frac{\pi}{2},\frac{\pi}{2} \right)\).

\emph{The only matching angle in the interval where the inverse trig
functions would be defined.}


\subsubsection{\texorpdfstring{ Inverse on the
inside}{ Inverse on the inside}}

What if we do the reverse?

\[\cos\left( \arccos(x) \right)\]

If we take a look at the arccosine function, it takes in a ratio between
\(\lbrack - 1,\ 1\rbrack\) and gives us an angle in
\(\lbrack 0,\ \pi\rbrack\). Then cosine takes an angle in its correct
angle domain according to how arccosine is defined. This means
\(x = \cos\left( \arccos(x) \right)\), ONLY IF
\(x \in \lbrack - 1,\ 1\rbrack\). \(\arccos(2) = \text{DNE}\)

Consider \(\sin\left( \arcsin(x) \right) = x\), as
\(\arcsin(x) = \theta \in \left\lbrack - \frac{\pi}{2},\frac{\pi}{2} \right\rbrack\),
ONLY IF \(x \in \lbrack - 1,\ 1\rbrack\). \(\arcsin(2) = \text{DNE}\);
don't try.

For tangent, we get the same conclusion.
\(\tan\left( \arctan(x) \right) = x\), which works for
\(x\mathbb{\in R}\). It is true always.

Therefore, having the inverse trigonometric function in the inside, you
will always get \(x\) as long as \(x\) is in the domain of the inverse
function.

We can conclude that:

\(\arccos\left( \cos(\theta) \right) = \theta\) for
\(\theta \in \lbrack 0,\ \pi\rbrack\)

\(\arcsin\left( \sin(\theta) \right) = \theta\) for
\(\theta \in \left\lbrack - \frac{\pi}{2},\frac{\pi}{2} \right\rbrack\)

Some questions:

\[\arccos\left( \cos\left( - \frac{\pi}{3} \right) \right) = \frac{\pi}{3}\]

(Work:
\(\cos\left( - \frac{\pi}{3} \right) = \frac{1}{2} = \cos\left( \frac{\pi}{3} \right)\);
\(\arccos\left( \cos\left( - \frac{\pi}{3} \right) \right) = \arccos\left( \frac{1}{2} \right) = \frac{\pi}{3} \neq - \frac{\pi}{3}\).
\(\frac{\pi}{3} \in \lbrack 0,\ \pi\rbrack\), required interval for
\(\arccos(\ldots)\) to be defined).

Now, how do we do this for any angle that is not in the special
triangles?


\subsection{Trig properties}

To easily calculate \(\arccos\left( \cos(\theta) \right)\) for any
\(\theta\), use the following properties:

\begin{itemize}
\item
  \(\cos(\theta) = \cos( - \theta)\) as cosine is an even function.
\item
  \(\cos(\theta \pm 2k\pi) = \cos(\theta)\ \forall k\mathbb{\in Z}\) as
  cosine is a periodic function. Note \(2\pi \approx 6.28\)
\end{itemize}

Sine is more complicated because the same properties do not hold.

\begin{itemize}
\item
  \(\sin( - \theta) = - \sin(\theta)\) as sine is an odd function -- you
  can move the negative sign in or out
\item
  \(\sin(\theta \pm \pi) = - \sin(\theta)\)

  \begin{itemize}
  \item
    As
    \(\sin(\theta)\cos(\pi) + \cos(\theta)\sin(\pi) = - \sin(\theta) + 0\)
    from the addition property.
  \item
    We can imply
    \(\sin\left( \theta \pm (2k + 1)\pi \right) = - \sin(\theta)\ \forall k\mathbb{\in Z}\)
  \end{itemize}
\item
  \(\sin(\theta + 2k\pi) = \sin(\theta)\ \forall k\mathbb{\in Z}\) as
  sine is a periodic function. Note \(2\pi \approx 6.28\)
\end{itemize}

Example:

\begin{align*}
\arcsin\left( \sin\left( \frac{5\pi}{6} \right) \right) &= - \arcsin\left( \sin\left( \frac{5\pi}{6} - \pi \right) \right) = \arcsin\left( - \sin\left( - \frac{\pi}{6} \right) \right) \\
&= \arcsin\left( \sin\left( \frac{\pi}{6} \right) \right) = \frac{\pi}{6}
\end{align*}



\subsection{Graphing inverse of itself for trig
functions}


\subsubsection{\texorpdfstring{ Sine}{ Sine}}

Sine is a periodic function. The arcsine of sine is also a periodic
function because sine is periodic. To figure out how
\(\arcsin\left( \sin(\theta) \right)\) will look like, you only need to
deduce what it will look like for \textbf{one period} of the function;
the function then repeats from there infinitely. By the way, the period
for \(\sin(\theta)\) is \(2\pi\).

Firstly, we know \(y = \arcsin\left( \sin(\theta) \right) = \theta\) if
\(\theta \in \left\lbrack - \frac{\pi}{2},\frac{\pi}{2} \right\rbrack\).
This gives us how the function behaves between that region, which covers
\(\frac{1}{2}\) of the period we need.

\textbf{Expanding regions}

Now, what is \(y\) when

\[\theta \in \left\lbrack \frac{\pi}{2},\frac{3\pi}{2} \right\rbrack\]

We can make the following deductions (REMEMBER: \(\theta\) STAYS
FAITHFUL TO ITS DEFINED INTERVAL,
\(\left\lbrack \frac{\pi}{2},\frac{3\pi}{2} \right\rbrack\)):

\[f(\theta) = \arcsin\left( \sin(\theta) \right) \neq 0\]

Use the property \(\sin(\theta \pm \pi) = - \sin(\theta)\), so subtract
\(\pi\) and add the negative sign to sine.

\[f(\theta) = \arcsin\left( - \sin(\theta - \pi) \right)\]

Then

\[\theta - \pi \in \left\lbrack - \frac{\pi}{2},\frac{\pi}{2} \right\rbrack\]

Unfortunately, arcsine and sine cannot cancel each other out. But since
we know

\[\theta - \pi \in \left\lbrack - \frac{\pi}{2},\frac{\pi}{2} \right\rbrack \Leftrightarrow \pi - \theta \in \left\lbrack - \frac{\pi}{2},\frac{\pi}{2} \right\rbrack\]

Then

\[f(\theta) = \arcsin\left( - \sin(\theta - \pi) \right) = \arcsin\left( \sin(\pi - \theta) \right) = \pi - \theta\]

Because
\(- \theta + \pi \in \left\lbrack - \frac{\pi}{2},\frac{\pi}{2} \right\rbrack\),
the above statement is true. This means

\[f(\theta) = - \theta + \pi,\ \theta \in \left\lbrack \frac{\pi}{2},\frac{3\pi}{2} \right\rbrack\]

By this point, we have found \(f(\theta)\) for
\(\theta \in \left\lbrack - \frac{\pi}{2},\frac{3\pi}{2} \right\rbrack\).

\textbf{Looping}

Note that \(\sin(\theta) = \sin(\theta + 2\theta)\). Therefore, we can
conclude our graph looks like this:

In other words: loop the following function:

\[f(\theta) = \left\{ \begin{matrix}
\theta,\ \theta \in \left\lbrack - \frac{\pi}{2},\frac{\pi}{2} \right\rbrack \\
 - \theta + \pi,\ \theta \in \left\lbrack \frac{\pi}{2},\frac{3\pi}{2} \right\rbrack \\
\end{matrix} \right.\ \]


\subsection{Converting between sine and
cosine}

Sine is just cosine transformed.

\[\sin(\theta) = \cos\left( \theta - \frac{\pi}{2} \right)\]

Trying to solve
\(\arccos\left( \sin(\theta) \right),\ \theta \in \left\lbrack \frac{\pi}{2},\frac{3\pi}{2} \right\rbrack\)

\begin{align*}
\arccos\left( \sin\theta \right) &= \arccos\left( \cos\left( \theta - \frac{\pi}{2} \right) \right),\ \theta - \frac{\pi}{2} \in \lbrack 0,\ \pi\rbrack \\
&= \theta - \pi/2
\end{align*}


As long as \(x\) is between the domain that lets cosine be injective,
\(\arccos\left( \cos(x) \right) = x\)

Now consider \(\theta\) being in a different region,
\(\theta \in \left\lbrack - \frac{\pi}{2},\frac{\pi}{2} \right\rbrack\)

Then
\(\arccos\left( \cos\left( \theta - \frac{\pi}{2} \right) \right),\ \theta - \frac{\pi}{2} \in \lbrack - \pi,\ 0\rbrack\).
Not in the same region where cos and arccosine would cancel each other
out. If you add \(2\pi\), it will overshoot, so don't. Consider negating
the input to cosine (taking advantage of the fact that cosine is an even
function).

\[= \arccos\left( \cos\left( - \theta + \frac{\pi}{2} \right) \right) = - \theta + \frac{\pi}{2} \in \lbrack 0,\ \pi\rbrack\ \text{(same\ as\ negating\ inequalities)}\]

Recall: placing the arc function inside the regular function will cancel
out, except the domain will be restricted. If you put the trig function
on the inside, any angle can be placed in.

\[\sec\left( \arccos(x) \right),\ x \in (0,\ 1\rbrack\]

\[\arccos(x) = \theta \in \lbrack 0,\ \pi\rbrack\]

\[\cos\left( \arccos(x) \right) = \cos(\theta) = x\]

By limiting \(x\) to be positive, it requires that \(\theta\) be
positive.

\[\frac{A}{H} = \frac{x}{1} = x = \cos(\theta)
\]

\[\sec\left( \arccos(x) \right)\]

But \(\arccos(x) = \theta\) so

\[\frac{1}{\cos\left( \arccos(x) \right)} = \frac{1}{x}\]

Now what about \(\sin\left( \arccos(x) \right)\) and knowing that \(x\)
is positive, then

\[\arccos(x) = \theta,\ \theta \in \left\lbrack 0,\frac{\pi}{2} \right\rbrack\ \text{(theta\ restriction\ from\ x\ forced\ positive)}\]

\[\sin(\theta) = \frac{l}{1} = \sqrt{1 - x^{2}}\]

This diagram: say that \(x = \cos(\theta)\). Now, calculate the height
of the opposite of \(\theta\). Using the Pythagoras theorem,
\(l = \sqrt{1 - x^{2}}\)


\section{Solving inverse functions}

Solving for an inverse function:

\begin{enumerate}
\def\labelenumi{\arabic{enumi}.}
\item
  Set \(y = f(x)\)
\item
  Isolate \(x\) on one side, \(x = g(y)\), where \(g(y)\) is the right
  side with only \(y\)
\item
  Write \(f^{- 1}(y) = g(y)\), then switch all \(y\) to \(x\), so
  \(f^{- 1}(x) = g(x)\). Note that switching the letter back to \(x\) is
  optional, but the moment you switch the \(x\) is NOT the same as the
  \(x\) of the original function.
\end{enumerate}


\subsection{Inverse Function Theorem}

Why do we use \(f^{- 1}\) to indicate inverse functions? Note:
\(\det\left( A^{- 1} \right) = \left( \det(A) \right)^{- 1} = \frac{1}{\det(A)}\).
The inverse of any number is its reciprocal. Now, get a function, and
let's say we have an inverse function: \(f\) and \(f^{- 1}\). Now, what
is \((f^{- 1})'\)? Note that you can interchange derivative operators
with inverse operators. This means:

\[\left( f^{- 1} \right)' = \left( f' \right)^{- 1}\]

\begin{align*}
f(x) &= y,\ x = f^{- 1}(y) \\
\left( f^{- 1} \right)'(y) &= \left( f'(x) \right)^{- 1} = \frac{1}{f'(x)} = \frac{1}{f'\left( f^{- 1}(y) \right)}
\end{align*}


Note: input changes when you swap a function's inverse and derivative
operator. You have a function in terms of \(y\), so substitute. The
input value \(y\) must show up on the right. \(x\) and \(y\) are related
by their inverse functions.


\subsection{Usage}

\[f(x) = \sin(x),\ f^{- 1}(x) = \arcsin(x)\]

Here, we are using \(x\) as both of the variables but we aren't supposed
to do that. If you are interested in taking the derivative of arcsine:

\begin{align*}
\left( f^{- 1} \right)'(x) &= \frac{1}{\cos\left( \arcsin(x) \right)} \\
\arcsin(x) &= \theta \\
x &= \sin(\theta) \\
l^{2} + x^{2} &= 1^{2} \\
l &= \sqrt{1 - x^{2}} = \cos(\theta)
\end{align*}




\section{Implicit differentiation}

Not functions

\begin{enumerate}
\def\labelenumi{\arabic{enumi}.}
\item
  Replace all \(y\) with \(y(x)\)
\item
  Take derivative against \(x\) on both sides.
\item
  Isolate \(y'(x) = \frac{\text{d}y}{\text{d}x}\)
\end{enumerate}

Remember to use the product and chain rule when taking derivatives.

For example:

\begin{align*}
\cos\left( xy \right) &= 1 + \sin(y) \\
\cos\left( x \cdot y(x) \right) &= 1 + \sin\left( y(x) \right) \\
\frac{d}{\text{d}x}\cos\left( x \cdot y(x) \right) &= \frac{d}{\text{d}x}\left\lbrack 1 + \sin\left( y(x) \right) \right\rbrack \\
- \sin\left( x \cdot y(x) \right) \cdot \left\lbrack 1 \cdot y(x) + x \cdot \frac{\text{d}y}{\text{d}x} \right\rbrack &= \cos\left( y(x) \right) \cdot \frac{\text{d}y}{\text{d}x}
\end{align*}


Isolate \(\frac{\text{d}y}{\text{d}x}\):

\begin{align*}
- \sin\left( x \cdot y(x) \right) \cdot y(x) - x \cdot \sin\left( x \cdot y(x) \right) \cdot \frac{\text{d}y}{\text{d}x} - \cos\left( y(x) \right) \cdot \frac{\text{d}y}{\text{d}x} &= 0 \\
\left\lbrack - x \cdot \sin\left( x \cdot y(x) \right) - \cos\left( y(x) \right) \right\rbrack \cdot \frac{\text{d}y}{\text{d}x} &= \sin\left( x \cdot y(x) \right) \cdot y(x) \\
\frac{\text{d}y}{\text{d}x} &= \frac{\sin\left( x \cdot y(x) \right) \cdot y(x)}{- x \cdot \sin\left( x \cdot y(x) \right) - \cos\left( y(x) \right)}
\end{align*}


Notation stuff:

\begin{align*}
y &= y(x) \\
y'(x) &= \frac{\text{d}y}{\text{d}x} = \frac{\text{d}}{\text{d}x}y(x) = y'
\end{align*}



\section{Related rates}

TAKE DERIVATIVE FIRST BEFORE PLUGGING IN VALUES

\begin{enumerate}
\def\labelenumi{\arabic{enumi}.}
\item
  Assign symbols to all variables in the problem
\item
  Writes an eqn that relates all the variables
\item
  Identify which variables change with time (If some quantity \(x\)
  changes as time passes, replace \(x\) with \(x(t)\)).
\item
  Take derivative against \(t\) on both sides using implicit
  differentiation
\item
  Plug in the values given for the variables in the problem.
\end{enumerate}

Example: vol of cone increases at 30\(m^{3}\) per minute. Diameter and
height are always equal. How fast is the height increasing when it is 10
meters tall?

\[v = \text{vol},\ r = \text{radius},\ h = ht\]

\[v = \frac{1}{3}\pi r^{2}h\]

Using assumption that \(d = h\),
\(d = 2r = h \Rightarrow r = \frac{h}{2}\)

\[v = \frac{1}{3}\pi\left( \frac{h}{2} \right)^{2}h = \frac{1}{12}\pi h^{3}\]

Because both change with time

\begin{align*}
v(t) &= \frac{1}{12}\pi{h(t)}^{3} \\
\frac{d}{\text{d}t}v(t) &= \frac{1}{12}\pi \cdot \frac{d}{\text{d}t}\left( h(t) \right)^{3} \\
\frac{\text{d}v}{\text{d}t} &= \frac{1}{12}\pi \cdot 3 \cdot h(t)^{2} \cdot \frac{\text{d}h}{\text{d}t} \\
30 &= \frac{1}{12} \cdot 3 \cdot \left( 10^{2} \right) \cdot \frac{\text{d}h}{\text{d}t} = 25\pi \cdot \frac{\text{d}h}{\text{d}t} \Rightarrow \frac{\text{d}h}{\text{d}t} = \frac{6}{5\pi}
\end{align*}


(Objective: find \(\frac{\text{d}h}{\text{d}t}\) or change in height
over time.)


The distance \(D\) is the distance between car W and car S. (W and S
indicate direction of car). \(x\) is the distance W driven so far, and
\(y\) is S driven so far. \(x(t),\ y(t)\) current position of car,
\(x'(t)\) and so on is the speed of the car right now.

\begin{align*}
x^{2} + y^{2} &= D^{2} \\
x(t)^{2} + y(t)^{2} &= D(t)^{2} \\
2x(t) \cdot x'(t) + 2y(t) \cdot y'(t) &= 2D(t) \cdot D'(t) \\
x(t) \cdot x'(t) + 2y(t) \cdot y'(t) &= D(t) \cdot D'(t)
\end{align*}


At
\(t = 2,\ x(t) = 50,\ y(t) = 120,\ D(t) = \sqrt{x^{2} + y^{2}} = \sqrt{50^{2} + 120^{2}} = \sqrt{2500 + 14400} = \sqrt{100} + \sqrt{25 + 144} = 10\sqrt{169} = 10 \cdot 13 = 130\)

\[50 \cdot 25 + 120 \cdot 60 = 1300D'(t)\]


\begin{align*}
x(t)^{2} + 200^{2} &= D(t)^{2} \\
2x(t)x'(t) &= 2D(t)D'(t) \\
x(t)x'(t) &= D(t)D'(t)
\end{align*}


When the police did the measurement, the radar indicates that the
distance between the police and the car is increasing by 70 km/h. That
is your \(D'(t)\)

Solve for \(x'(t) \approx 100\)

\begin{align*}
200x'(t) &= 100\sqrt{8} \cdot 70 \\
x'(t) &= \frac{100\sqrt{8} \cdot 70}{200} = \frac{1}{2}\sqrt{8} \cdot 70 = 35\sqrt{8}
\end{align*}



\section{MVT revisited}

Consequences of the MVT

\[f(b) = f(a) \Rightarrow \exists c \in (a,\ b),\ f'(c) = 0\]

Constant function

If \(f'(x) = 0\forall x\) in an interval, \(\Rightarrow f(x)\) is a
constant function.

If \(f'(x) \neq 0\forall x\) in an interval, then \(f(x)\) is 1-1.

If \(f'(x) > 0\forall x\) in an interval, then \(f(x)\) is
increasing.

If \(f'(x) < 0\) \ldots{} \(f(x)\) is decreasing.

Local max: A point \(c\) is a local max if \(\exists\) interval
\((c - l,\ c + l)\) such that for all \(x\) in interval,
\(f(x) \leq f(c)\) when \(f(x)\) is defined.

Local extrema: in the interior of interval \(\lbrack a,\ b\rbrack\) then
\(c\) is a critical point \(f'(c) = 0\) or \(f'(c)\) DNE.

\begin{align*}
f(x) &= \arcsin\left( \frac{1 - x}{1 + x} \right) + 2\arctan\left( \sqrt{x} \right) \\
&= \frac{1}{\sqrt{1 - \left( \frac{1 - x}{1 + x} \right)^{2}}} \cdot \left( - \frac{2}{(1 + x)^{2}} \right) \\
&= - \frac{2}{\sqrt{4x}(1 + x)} + \frac{1}{(1 + x)\sqrt{x}} = - \frac{1}{\sqrt{x}(1 + x)} + \frac{1}{(1 + x)\sqrt{x}} = 0
\end{align*}


Note:

\[f(0) = c = \arcsin\left( \frac{1}{1} \right) + 2\arctan(0) = \frac{\pi}{2} + 0\]

(Test a value somewhere in the interval.)

Deducing domain of the function:

\[\arcsin\left( \frac{1 - x}{1 + x} \right) + 2\arctan\left( \sqrt{x} \right)
\]

True

\begin{align*}
1 - x &\leq 1 + x \\
1 &\leq 1 + 2x \\
0 &\leq 2x
\end{align*}


Another question

Find all functions defined on \(\mathbb{R}\) that has \(f'(x) = 3\)
everywhere.

That is \(f(x) = 3x + C\), where \(C\) is a constant. We can mention it
in another way: there is an infinite number of functions that have
\(f'(x) = 3\) everywhere in the form \(f(x) = 3x + C\).

Suppose \(g(x) \neq 3x + C\) and \(g'(x) = 3\) everywhere. How do we
show that? The trick here is to consider:

\[f'(x) - g'(x) = 0,\ as\ f'(3) - g'(3) = 0\]

This means

\[f - g = D,\ \text{a\ constant}\]

Is a constant function.

However, if you solve for \(g\):
\(f - g = D \Rightarrow g = f - D \Rightarrow g = 3x + C - D\)

But \(C - D\ \)is just another constant.


\subsection{At most some zeroes}

Suppose \(f'(x)\) has only \(k\) zeroes on \(\lbrack a,\ b\rbrack\).
Then \(f(x)\) can have at most \(k + 1\) zeroes. This can be proved by
induction.


\section{Maximum and minimum (second derivative
test)}

What to do at each critical point:

\begin{itemize}
\item
  \(f''(a) > 0\), \(a\) is a local minimum.
\item
  \(f''(a) < 0\), then local max
\item
  \(f''(a) = 0\) then it is unclear
\item
  We can also use the first derivative test for the local max or min.
\end{itemize}

Compare all values of \(f(a)\) for all critical points and endpoints of
an interval. The largest value is a global maximum. The smallest value
is an absolute (global) minimum.


\subsection{Absolute max and min}

For
\(\lim_{x \rightarrow \infty}{f(x)} = + \infty = \lim_{x \rightarrow - \infty}{f(x)}\)
then there exists an absolute minimum.

For
\(\lim_{x \rightarrow \infty}{f(x)} = - \infty = \lim_{x \rightarrow - \infty}{f(x)}\)
then there exists an absolute max.


For
\(\lim_{x \rightarrow \infty}{f(x)} = \pm \infty,\ \lim_{x \rightarrow - \infty}{f(x)} = \mp \infty\)
then no absolute max or min.


\subsubsection{\texorpdfstring{
Optimization}{ Optimization}}

\begin{enumerate}
\def\labelenumi{\arabic{enumi}.}
\item
  Identify quantity to be maximized/minimized. Set as \(y\).
\item
  Identify changeable quantity. Set as \(x\).
\item
  Write an equation that relates all variables and draw a diagram if
  needed.
\item
  Calculate \(\frac{\text{d}y}{\text{d}x}\) and set to 0.
\item
  Solve for \(x\) and the quantities needed.
\item
  Context helps you determine which \(x\) is correct.
\item
  -\\
\end{enumerate}


\section{Concavity and Point of
Inflection}

If \(f'(x) > 0\) function increasing, negative then the function is
decreasing.

For \(f''(x)\), positive \(\Rightarrow\) concave up; negative then
concave down.

A point of inflexion is when \(f''(x) = 0\) and the sign changes. If
\(f''(x)\) changes sign at \(x\) but \(f''(x)\) DNE it can still
be a point of inflexion. Test points around the point of inflexion.


\section{Graph sketching}

Draw a decent enough picture so calculus can be done on it. We need the
following:

\begin{itemize}
\item
  Domain
\item
  Vertical asymptotes: \(\lim_{x \rightarrow a}{f(x)} = \pm \infty\)
\item
  Horizontal asymptotes and slant asymptotes (when the power on top is
  greater than the power on the bottom). Obtain slant asymptotes using
  long division.
\item
  Points of intersection with axis (\(x\) and \(y\) intercept)
\item
  Critical points. Set derivative equal to 0.
\item
  Points of inflexion, but don't do that.
\item
  Others.
\end{itemize}


\subsection{Examples}

\[\frac{x^{3}}{x^{2} + 1}\]

No VA, as \(x^{2} \neq - 1\)

May have slant asymptotes. \(\lim_{x \rightarrow \infty}f = + \infty\),
\(\lim_{x \rightarrow - \infty}f = - \infty\)

Occurs when power on top of the rational function is higher than the
bottom.


The remainder means if I were to remove the remainder from the
beginning, then the division would go through. This means if you
subtract the remainder from the numerator:

\[\frac{x^{3} - ( - x)}{x^{2} + 1} = x\]

You can have this:

\[\frac{x^{3}}{x^{2} + 1} = x + \left( \frac{- x}{x^{2} + 1} \right)\]

The function highlighted in orange will always have a lower degree on
the top.

If we now take
\(\lim_{x \rightarrow \pm \infty}\left( - \frac{x}{x^{2} + 1} \right)\),
we will get 0, as the degree is larger on the bottom. This means, in the
long run, you will end with \(x\). This means

\[{\lim_{x \rightarrow \pm \infty}\left( \frac{x^{3}}{x^{2} + 1} \right)} \approx x\]

This function, \(y = x\) is power 1 is because the numerator \(x^{3}\)
and the denominator \(x^{2} + 1\) differ by one power. When the top
power is 1 more than the bottom, the quotient will be a linear function.
The function essentially becomes linear as \(x\) approaches infinity. In
general, if you have \(\frac{x^{4}}{x^{2} + 1}\), you will get a
quadratic as the quotient, so instead of having a straight line to what
the function will conform to, you will get a parabola instead.

Point of intersect: \((0,\ 0)\) (just plug it in)

Take the derivative and find critical points:

\begin{align*}
f' &= \frac{3x^{2}\left( x^{2} + 1 \right) - 2x \cdot x^{3}}{\left( x^{2} + 1 \right)^{2}} = 0 \\
3x^{4} + 3x^{2} - 2x^{4} &= 0 \\
x^{2}\left( x^{2} + 3 \right) &= 0 \\
x^{2} &= 0,\ x = 0
\end{align*}


At \(x = 0\), this is the only critical point.

Order of zeroes: The only zero is at \(x = 0\), which is what we found.
If we take a look at the function, \(f = \frac{x^{3}}{x^{2} + 1}\). This
means that this is an order 3, as the power on \(x\) is 3. You only need
to worry about the numerator, is that the idea is, when we are
interested in order of 0, we only care how the function behaves near the
zero, which turns out to be \(x = 0\). When \(x\) is near 0, the
denominator \(x^{2} + 1\) basically becomes a constant, approaching 1.
Near \(x = 0\), the function is essentially \(\frac{x^{3}}{1}\). That is
the reason why \(x^{3}\) governs the behavior of the function.


\subsection{Order of zeroes}

Aside:

\[g = \frac{(x - 1)^{2}(x + 2)^{3}}{x - 2}\]

Zero: \(x = 1,\ x = - 2\). Order of \(x = 1\) is 2; order of \(x = - 2\)
is 3. If the order is even, the graph will bounce off. If the order is
odd, the function will make a cubic-like wiggle across the \(x\)-axis.


\section{L'Hopital's rule}

L'Hopital's rule makes no sense of why it works until you get to Taylor
series, but not every function can be written as a Taylor series.

When

\[\frac{\lim{f(x)}}{\lim{g(x)}} = \frac{0}{0}\ \text{or} \pm \frac{\infty}{\infty}\]

Then

\[\lim\left( \frac{f(x)}{g(x)} \right) = \lim\left( \frac{f'(x)}{g'(x)} \right)\]


\subsection{Examples}

\[\lim_{x \rightarrow 0}\frac{x^{3}}{x^{2} + x} = \frac{0}{0}\]

\[= \lim_{x \rightarrow 0}\frac{3x^{2}}{2x + 1} = \frac{0}{1} = 0\]

This is done by taking the derivative of the numerator and the
denominator.


\subsection{Flip and multiply}

Case of \(0 \cdot \infty\)

\[\lim\left( f(x) \cdot g(x) \right) = 0 \cdot \infty\]

Turn it into a fraction by forcing either \(f\) or \(g\) onto the
denominator. Choose either, better the simpler one. Whatever gets
flipped turns from \(\pm \infty\) to 0 or the other way around.

\begin{align*}
\lim\left( \frac{f(x)}{\frac{1}{g(x)}} \right) &= \frac{0}{0}\text{\ or\ } \pm \frac{\infty}{\infty} \\
\lim\left( \frac{g(x)}{\frac{1}{f(x)}} \right) &= \frac{0}{0}\text{\ or\ } \pm \frac{\infty}{\infty}\
\end{align*}


Then proceed with L'Hopital's rule.


\subsection{Non-indeterminate forms}

\[\frac{0}{\infty} = 0\]

\[\frac{\infty}{0} = \infty\]


\subsection{How would we have figured that
out?}

Suppose \(f(x) = a_{0} + a_{1}x + a_{2}x^{2} + a_{3}x^{3} + \ldots\) (an
infinite polynomial) and
\(g(x) = b_{0} + b_{1}x + b_{2}x^{2} + b_{3}x^{3} + \ldots\) and have

\[\lim_{x \rightarrow 0}\frac{f(x)}{g(x)} = \frac{a_{0} + a_{1}x + a_{2}x^{2} + a_{3}x^{3} + \ldots}{b_{0} + b_{1}x + b_{2}x^{2} + b_{3}x^{3} + \ldots}\]

Since every term at the back contains an \(x\), most of the time, we
should end up with

\[\frac{a_{0}}{b_{0}}\]

L'Hopital's rule says that this answer is true most of the time, but if
\textbf{both} \(b_{0} = 0\ \text{and}\ a_{0} = 0\), then we end up with
an indeterminate form.

When you take L'Hopital's rule, only these have an impact:

\[\frac{a_{0} + a_{1}x + a_{2}x^{2} + a_{3}x^{3} + \ldots}{b_{0} + b_{1}x + b_{2}x^{2} + b_{3}x^{3} + \ldots}\]

With \(a_{0}\) and \(b_{1}\) disappearing, we should actually get

\[\frac{a_{1} + a_{2}x + a_{3}x^{2} + \ldots}{b_{1} + b_{2}x + b_{3}x^{2} + \ldots} = \frac{a_{1}}{b_{1}}\]

Repeat process if \(a_{1}\) and \(b_{1}\) are zero.

When you take the derivative of \(f(x)\) and \(g(x)\), you remove
\(a_{0}\) and \(b_{0}\). Meaning:

\[\frac{f'(x)}{g'(x)} = \frac{a_{1} + a_{2}x + a_{3}x^{2} + \ldots}{b_{1} + b_{2}x + b_{3}x^{2} + \ldots}\]

Note that this fraction is completely different from \(\frac{f}{g}\).
This gets rid of \(a_{0}\) and \(b_{0}\).

Now we have \(\frac{a_{1}}{b_{1}}\). This is a way to extract it.

Analytic functions can be written as polynomials, and nearly all
functions aren't. Yet, the functions we work with the most often are
analytic functions.


\subsection{Exponent cases}

When the base of an exponent \textbf{includes} \(x\), then use the
formula:

\[f(x)^{g(x)} = \left( e^{\ln\left( f(x) \right)} \right)^{g(x)} = e^{g(x) \cdot \ln\left( f(x) \right)}\]

\begin{align*}
f(x) &= e^{\ln\left( f(x) \right)} \\
f(x)^{g(x)} &= \left( e^{\ln\left( f(x) \right)} \right)^{g(x)} \\
x^{x} &= e^{\ln\left( x^{x} \right)} = e^{x\ln(x)}
\end{align*}


For example:

\[\lim_{x \rightarrow 0^{+}}\left( \tan(2x) \right)^{x} = \lim_{x \rightarrow 0^{+}}e^{x \cdot \ln\left( \tan(2x) \right)}\]

Using the property of limits, we can move the limit onto the exponent.

\[e^{\lim_{x \rightarrow 0^{+}}{x \cdot \ln\left( \tan(2x) \right)}}\]

So we only need to compute:


Simplify stuff the moment you take the derivative.


You may have to take L'Hopital's again.

\end{document}
